\documentclass{article}
%packages
\usepackage{fullpage}
\usepackage{xspace}
\usepackage{graphicx,color}
\usepackage{hyperref}
\usepackage{graphicx}
\usepackage[dvipsnames]{xcolor}
\usepackage{amsmath, amsthm, amscd, amsfonts, amssymb}
\usepackage{graphicx}
\usepackage{stmaryrd}
\usepackage{quiver}
\usepackage{float}
\usepackage{mathtools}
\usepackage{mathrsfs}
\usepackage{multicol}
\usepackage{tikz-cd}
\usepackage[shortlabels]{enumitem}
\usepackage{stackrel}
\usepackage{cancel}
\usepackage[framemethod = TikZ]{mdframed}

%theorems
\newtheorem{theorem}{Teorema}[subsection]

\newtheorem{algorithm}{Algorithm}
\newtheorem{axiom}{Axioma}
\newtheorem{case}{Case}
\newtheorem{claim}{Afirmación}
\newtheorem{conclusion}{Conclusión}
\newtheorem{condition}{Condition}
\newtheorem{conjecture}[theorem]{Conjetura}
\newtheorem{corollary}[theorem]{Corolario}
\newtheorem{criterion}{Criterion}
\newtheorem{proposition}[theorem]{Proposición}
\newtheorem{lemma}[theorem]{Lema}

\theoremstyle{definition}
\newtheorem{definition}[theorem]{Definición}

\newtheorem{exercise}{Ejercicio}

\newtheorem*{notation}{Notación}
\newtheorem{problem}{Problema}[section]
\newtheorem*{obs}{Observación}
\newtheorem*{addendum}{Addendum}

\newtheorem*{tarea}{\wea{Tarea}}

\theoremstyle{remark}
\newtheorem{remark}[theorem]{Remark}
\newtheorem*{solution}{\textbf{Solución}}
\newtheorem{summary}{Summary}
\newtheorem{example}[theorem]{\textbf{Ejemplo}}

\numberwithin{equation}{section}

\renewcommand*{\proofname}{\textbf{Demostración}}

\surroundwithmdframed[
  topline=false,
  rightline=false,
  bottomline=false,
  leftmargin=\parindent,
  skipabove=\medskipamount,
  skipbelow=\medskipamount
]{proof}

\surroundwithmdframed[
  topline=false,
  rightline=false,
  bottomline=false,
  leftmargin=\parindent,
  skipabove=\medskipamount,
  skipbelow=\medskipamount
]{solution}

%commands

%Main letters
\newcommand{\A}{\mathbb{A}}
\newcommand{\B}{\mathbb{B}}
\newcommand{\C}{\mathbb{C}}
\newcommand{\D}{\mathbb{D}}
\newcommand{\F}{\mathbb{F}}
\newcommand{\N}{\mathbb{N}}
\renewcommand{\P}{\mathbb{P}}
\newcommand{\Q}{\mathbb{Q}}
\newcommand{\R}{\mathbb{R}}
\newcommand{\Z}{\mathbb{Z}}

\newcommand{\Acal}{\mathcal{A}}
\newcommand{\Bcal}{\mathcal{B}}
\newcommand{\Ccal}{\mathcal{C}}
\newcommand{\Dcal}{\mathcal{D}}
\newcommand{\Ecal}{\mathcal{E}}
\newcommand{\Fcal}{\mathcal{F}}
\newcommand{\Gcal}{\mathcal{G}}
\newcommand{\Hcal}{\mathcal{H}}
\newcommand{\Ical}{\mathcal{I}}
\newcommand{\Jcal}{\mathcal{J}}
\newcommand{\Kcal}{\mathcal{K}}
\newcommand{\Lcal}{\mathcal{L}}
\newcommand{\Mcal}{\mathcal{M}}
\newcommand{\Ncal}{\mathcal{N}}
\newcommand{\Ocal}{\mathcal{O}}
\newcommand{\Pcal}{\mathcal{P}}
\newcommand{\Qcal}{\mathcal{Q}}
\newcommand{\Rcal}{\mathcal{R}}
\newcommand{\Scal}{\mathcal{S}}
\newcommand{\Tcal}{\mathcal{T}}
\newcommand{\Ucal}{\mathcal{U}}
\newcommand{\Vcal}{\mathcal{V}}
\newcommand{\Wcal}{\mathcal{W}}
\newcommand{\Xcal}{\mathcal{X}}
\newcommand{\Ycal}{\mathcal{Y}}
\newcommand{\Zcal}{\mathcal{Z}}

\newcommand{\lam}{\lambda}
\newcommand{\e}{\varepsilon}
\newcommand{\w}{\omega}
\renewcommand{\a}{\alpha}

%symbols
\newcommand{\vphi}{\varphi}
\newcommand{\oo}{\infty}
\newcommand{\longto}{\longrightarrow}
\newcommand{\embed}{\hookrightarrow}
\newcommand{\acts}{\curvearrowright}
\newcommand{\actsby}[1]{\stackrel{#1}{\acts}}
\renewcommand{\c}{\cdot}
\renewcommand{\bar}{\overline}
\newcommand{\normal}{\trianglelefteqslant}
\renewcommand{\tilde}{\widetilde}
\renewcommand{\hat}{\widehat}

%special functions or sets
\newcommand{\sgn}{\mathrm{sgn}}
\newcommand{\id}{\mathrm{id}}
\newcommand{\im}{\mathrm{im}\,}
\newcommand{\op}{\mathrm{op}}
\newcommand{\orb}{\mathrm{orb}}
\newcommand{\GL}{\mathrm{GL}}
\newcommand{\SL}{\mathrm{SL}}
\newcommand{\Sym}{\mathrm{Sym}}
\newcommand{\Aut}{\mathrm{Aut}}
\newcommand{\Inn}{\mathrm{Inn}}
\newcommand{\Hom}{\mathrm{Hom}}
\newcommand{\stab}{\mathrm{stab}}
\newcommand{\Cent}{\mathrm{Cent}}
\newcommand{\ab}{\mathrm{ab}}
\newcommand{\sop}{\mathrm{sop}}
\newcommand{\Cay}{\mathrm{Cay}}
\newcommand{\ord}{\mathrm{ord}}
\newcommand{\lcm}{\mathrm{lcm}}
\newcommand{\Fix}{\mathrm{Fix}}
\newcommand{\Norm}{\mathrm{Norm}}

%different bracket types
\newcommand{\abs}[1]{\left| #1 \right|}
\newcommand{\norm}[1]{\left\lVert #1 \right\rVert}
\newcommand{\br}[1]{\left( #1 \right)}
\newcommand{\set}[1]{\left\{ #1 \right\}}
\newcommand{\cor}[1]{\left[ #1 \right]}
\newcommand{\gen}[1]{\left\langle #1 \right\rangle}
\newcommand{\matr}[1]{\begin{pmatrix}#1\end{pmatrix}}

\newcommand{\znz}[1]{\frac{\Z}{#1 \Z}}
\newcommand{\zmz}[1]{\Z/ #1 \Z}

%If I mess up while typing
\newcommand{\rac}[2]{\frac{#1}{#2}}
\newcommand{\olon}{\colon}
\newcommand{\ubseteq}{\subseteq}
\newcommand{\ubset}{\subset}
\newcommand{\irc}{\circ}
\renewcommand{\o}{\to}

\newcommand{\ds}{\displaystyle}

\renewcommand*\contentsname{Contenidos}

\newcommand{\wea}[1]{\textcolor{BlueViolet}{\textbf{\emph{#1}}}}
\newcommand{\fecha}[1]{\begin{flushright} \underline{#1} \end{flushright}}
\newcommand{\dem}[1]{\textbf{Demostración (#1)}}

\makeatletter
\providecommand*{\cupdot}{%
  \mathbin{%
    \mathpalette\@cupdot{}%
  }%
}
\newcommand*{\@cupdot}[2]{%
  \ooalign{%
    $\m@th#1\cup$\cr
    \sbox0{$#1\cup$}%
    \dimen@=\ht0 %
    \sbox0{$\m@th#1\cdot$}%
    \advance\dimen@ by -\ht0 %
    \dimen@=.5\dimen@
    \hidewidth\raise\dimen@\box0\hidewidth
  }%
}

\providecommand*{\bigcupdot}{%
  \mathop{%
    \vphantom{\bigcup}%
    \mathpalette\@bigcupdot{}%
  }%
}
\newcommand*{\@bigcupdot}[2]{%
  \ooalign{%
    $\m@th#1\bigcup$\cr
    \sbox0{$#1\bigcup$}%
    \dimen@=\ht0 %
    \advance\dimen@ by -\dp0 %
    \sbox0{\scalebox{2}{$\m@th#1\cdot$}}%
    \advance\dimen@ by -\ht0 %
    \dimen@=.5\dimen@
    \hidewidth\raise\dimen@\box0\hidewidth
  }%
}
\makeatother

%title
\title{MPG3200 -- Álgebra I}
\author{Felipe Inostroza \\ Prof. Sebastián Barbieri}
\date{Primer Semestre 2026}

\begin{document}

\maketitle
\tableofcontents

\newpage
\fecha{Clase 1. Lunes 9 de Marzo}
\section{Grupos}

\begin{definition}
  Un \wea{grupo} es un par $(G,\mathrm{op})$ donde $G$ es un conjunto y $\mathrm{op} \colon G \times G \to G$ tal que si denotamos $g\cdot h = \mathrm{op}(g,h)$, $\forall g,h \in G$, entonces
  \begin{enumerate}[(1)]
    \item $\forall x,y,z \in G$, $(x \cdot y) \cdot z = x \cdot (y \cdot z)$;
    \item $\exists e \in G$ tal que $\forall x \in G$, $e \cdot x = x \cdot e = x$.
    \item $\forall x \in G$, $\exists y \in G$ tal que $x \cdot y = y \cdot x = e$.
  \end{enumerate}
\end{definition}

\begin{definition}
  $(G,\op)$ se dice \wea{abeliano} si $\op(x,y) = \op(y,x)$, $\forall x,y \in G$.
\end{definition}

\begin{obs}
  En todo grupo, el neutro es único.
\end{obs}

\begin{notation}
  El neutro de un grupo $(G,\cdot)$ se denota:
  \begin{itemize}
    \item $e, e_G$,
    \item $1, 1_G$,
    \item $0, 0_G$ (exclusivamente si $G$ es abeliano)
    \item $\id$.
  \end{itemize}
  La operación se suele denotar por
  \begin{itemize}
    \item $\op(x,y) = x \cdot y$ (en general para grupos no abelianos)
    \item $\op(x,y) = x+y$ (exclusivamente para grupos abelianos)
  \end{itemize}
  Si $xy = yx = e$, escribimos $y = x^{-1}$.
  Si $x+y=y+x=e$, escribimos $y = -x$.
\end{notation}

\begin{definition}
  Sea $G$ grupo. Un subconjunto $H \subset G$ se dice \wea{subgrupo} (denotado $H \leqslant G$) si $H$ con la operación heredada de $G$ es grupo.
\end{definition}

\begin{proposition}
  Sea $H \subset G$ no vacío. Entonces $H \leqslant G \iff \forall h,h' \in H, h'h^{-1} \in H$.
\end{proposition}

\begin{example}
  $(\Z,+) \leqslant (\Q,+) \leqslant (\R,+) \leqslant (\C,\leqslant)$.
\end{example}

\begin{example}
  Si $(V,+,\cdot)$ es un espacio vectorial, entonces $(V,+)$ es un grupo abeliano. Por ejemplo $(\R^n,+)$ y $(\R[x],+)$ son grupos abelianos.
\end{example}

\begin{example}
  Si $M_n(\R)$ es el conjunto de matrices de $n \times n$ sobre $\R$, entonces
  \[\GL_n(\R) = \{A \in M_n(\R) \colon \det(A) \neq 0\}\]
  es un grupo no abeliano.
\end{example}

\begin{example}
  $\SL_n(\R) = \{A \in \GL_n(\R) \colon \det(A) = 1\}$ es un subgrupo de $\GL_n(\R)$.
\end{example}

\begin{example}
  Dado un conjunto $X$, consideramos el \wea{grupo simétrico}
  \[\Sym(X) = \{f \colon X \to X \text{ biyectivas}\}.\]
  Tenemos que $(\Sym(X),\circ)$ es un grupo.
\end{example}

\begin{example}
  $S^1 = \{z \in \C \colon |z| = 1\}$ con la multiplicación.

  Fijemos $n \geq 1$ entero.
  \[\Z_n = \{z \in S^1 \colon z^n = 1\}.\]
  Se tiene que $\Z_n \leqslant S^1$.
\end{example}

\begin{definition}
  Sean $G,H$ grupos. Un \wea{homomorfismo} es una función $$\phi \colon G \to H$$ tal que $\forall x, y \in G$, $$\phi(x\c y) = \phi(x) \c \phi(y).$$
\end{definition}

\begin{obs}
  $\phi(1_G) = 1_H$, $\phi(g^{-1}) = (\phi(g))^{-1}$.
\end{obs}

\begin{definition}
  Sea $\phi \colon G \to H$ homomorfismo de grupos. La \wea{imagen} y el \wea{kernel} de $\phi$ son (resp.)
  \begin{align*}
    \im(\phi) &= \{h \in H \colon \exists g \in G, \phi(g) = h\}, \\
    \ker(\phi) &= \{g \in G \colon \phi(g) = 1_H\}.
  \end{align*}
\end{definition}

\begin{proposition}
  $\im(\phi) \leqslant H$, $\ker(\phi) \leqslant G$.
\end{proposition}
\begin{proof}
  \
  \begin{itemize}
    \item \boxed{\im\phi} Primero notemos que $\phi(e_G) = e_H \in \im(\phi)$, por lo que $\im(\phi) \neq \varnothing$. Ahora, sean $h_1,h_2 \in \im(\phi)$, entonces existen $g_1,g_2 \in G$ tales que $h_i = \phi(g_i)$. Luego $$h_1 h_2^{-1} = \phi(g_1) (\phi(g_2))^{-1} = \phi(g_1) \phi(g_2^{-1}) = \phi(g_1 g_2^{-1}) \in \im(\phi)$$ y como $h_1,h_2$ son arbitrarios, tenemos que $\im(\phi) \leqslant H$.
    \item \boxed{\ker(\phi)} Sabemos que $\phi(e_G) = e_H$, por lo que $e_G \in \ker(\phi)$, y entonces $\ker(\phi) \neq \varnothing$. Sean $g,h \in \ker(\phi)$, entonces $\phi(g) = \phi(h) = e_H$. Luego $$\phi(gh^{-1}) = \phi(g) (\phi(h))^{-1} = e_H \c e_H = e_H$$ y entonces $gh^{-1} \in \ker(\phi)$. Por lo tanto $\ker(\phi) \leqslant G$.
  \end{itemize}
\end{proof}

\begin{exercise}
  Un homomorfismo $\phi \colon G \to H$ es inyectivo si y solo si $\abs{\ker(\phi)} = 1$.
\end{exercise}
\begin{proof}
  \boxed{\implies}
  Supongamos que $\phi \colon G \to H$ es un homomorfismo inyectivo. Sea $x \in \ker(\phi)$, entonces $\phi(x) = e_G = \phi(e_G)$, y como $\phi$ es inyectivo, tenemos que $x = e_G$. Por lo tanto $\ker(\phi) = \{e_G\}$, y entonces $\abs{\ker(\phi)} = 1$.

  \boxed{\impliedby}
  Supongamos que $\abs{\ker(\phi)} = 1$. Sean $g,h \in G$ tales que $\phi(g) = \phi(h)$. Entonces
  $$e_H = \phi(g) \c (\phi(h))^{-1} = \phi(gh^{-1})$$
  por lo que $gh^{-1} \in \ker(\phi)$. Como $e_G \in \ker(\phi)$, y $|\ker(\phi)| = 1$, tenemos que $\ker(\phi) = \{e_G\} \ni gh^{-1}$. Luego $gh^{-1} = e_G \implies g=h$. Por lo tanto $\phi$ es inyectivo.
\end{proof}

\begin{definition}
  El \wea{grupo trivial} es un grupo con un único elemento. Se denota por $G = 1$.
\end{definition}

\begin{example}
  Sea $n \geq 1$ entero, y consideremos
  \begin{align*}
    \psi_n \colon S^1 &\to S^1 \\
    z &\mapsto z^n
  \end{align*}
  Esta función es homomorfismo, con $\im(\psi_n) = S^1$, $\ker(\psi_n) = \Z_n$.
\end{example}

\begin{example}
  Toda transformación lineal $T \colon \R^n \to \R^n$ es homomorfismo.
  \begin{align*}
    \phi \colon \Z &\to \GL_2(\Z) \\
    n &\mapsto \matr{1&n\\0&1}
  \end{align*}
\end{example}

\begin{notation}
  Sea $\phi \colon G \o H$ homomorfismo. Diremos que
  \begin{itemize}
    \item $\phi$ es \wea{monomorfismo} si $\phi$ es inyectivo.
    \item $\phi$ es \wea{epimorfismo} si $\phi$ es sobreyectivo.
    \item $\phi$ es \wea{isomorfismo} si $\phi$ es biyectivo.
    \item $\phi$ es \wea{endomorfismo} si $G=H$.
    \item $\phi$ es \wea{automorfismo} si $\phi$ es biyectiva y $G=H$.
  \end{itemize}
  Si existe un isomorfismo entre $G$ y $H$, decimos que son \wea{isomorfos}, denotado por $G \simeq H$.
\end{notation}

%% demostrar
\begin{example}
  Sea $K = \set{\matr{x&-y\\y&x} \in M_2(\R)}$ con la suma. Tenemos que $(K,+) \simeq (\C,+)$.

  Si $K^\ast = K \setminus \set{\matr{0&0\\0&0}}$, entonces $(K^\ast,\c) \simeq (\C \setminus \{0\}, \c)$.
\end{example}

\begin{example}
  $(\zmz n, +) \simeq (\Z_n, \cdot)$.
\end{example}

\begin{exercise}
  Entre los siguientes grupos, determine cuales son isomorfos entre si y cuales no.
  \[(\Z,+), \quad (\Q,+), \quad (\R,+), \quad (\C,+).\]
\end{exercise}
\begin{proof}\
  \begin{itemize}
    \item \boxed{\Z \not\simeq \Q} Supongamos que existe un isomorfismo $\phi \colon \Z \to \Q$. Sea $q = \phi(1)$, luego existe $u \in \Z$ tal que $\phi(u) = q/2$. Entonces $$\phi(u+u) = \phi(u) + \phi(u) = q$$ y entonces $u+u = 1$. $\to\gets$
    \item \boxed{\Z \not\simeq \R} Recordemos que $|\Z| \neq |\R|$, es decir, no existe biyección entre $\Z$ y $\R$, y por lo tanto no puede existir un isomorfismo $\Z \to \R$.
    \item \boxed{\Z \not\simeq \C} Mismo argumento.
    \item \boxed{\Q \not\simeq \R} Mismo argumento.
    \item \boxed{\Q \not\simeq \C} Mismo argumento.
    \item \boxed{\R \simeq \C} Notemos que $\R$ y $\C$ son espacios vectoriales sobre $\Q$, y ambos tienen dimensión $\aleph_0$ (con base de Hamel), por lo que existe un isomorfismo de espacios vectoriales $\phi \colon \R \to \C$. La restricción de este isomorfismo con la suma nos da un isomorfismo de grupos.
  \end{itemize}
\end{proof}

\begin{definition}
  Sea $G$ un grupo y $g \in G$. La \wea{conjugación} es la aplicación
  \begin{align*}
    C_g \colon G &\to G \\
    h &\mapsto ghg^{-1}.
  \end{align*}
  $C_g$ es un automorfismo.
\end{definition}

\begin{notation}
  Dado un grupo $G$,
  $$\Aut(G) = \{\varphi \colon G \to G \text{ automorfismo}\}$$
  con la composición es el \wea{grupo de automorfismos de $G$}.
  $$\Inn(G) = \{\varphi \colon G \to G \mid \exists g \in G, \varphi = C_g\}$$
  con la composición es el \wea{grupo de automorfismos internos}.
\end{notation}

\begin{obs}
  $\Inn(G) \leqslant \Aut(G)$.
\end{obs}

\begin{example}
  Si $(G,+)$ es abeliano, entonces $\Inn(G) \simeq 1$.
\end{example}

\begin{exercise}
  Pruebe que $$\Aut(\Z^d,+) \simeq (\GL_d(\Z), \c).$$
\end{exercise}

%% demostrar cosas

\begin{exercise}
  Mostrar que $(\Q,+) \not\simeq (\Q^2,+)$ y $(\Z,+) \not\simeq (\Z^2,+)$
\end{exercise}

\newpage
\fecha{Clase 2. Miércoles 11 de Marzo}

\begin{proposition}
  Dados $X$ e $Y$ conjuntos, entonces
  \[\Sym(X) \simeq \Sym(Y) \iff |X| = |Y|\]
\end{proposition}
\begin{proof}
  \
  \begin{itemize}
    \item \boxed{\impliedby} Supongamos que $|X| = |Y|$. Entonces existe $f \colon X \to Y$ biyectiva. Definimos $$\phi \colon \Sym(X) \to \Sym(Y)$$ dada por $\phi(\sigma) = f \circ \sigma \circ f^{-1}$ para todo $\sigma \in \Sym(X)$.
    \item \boxed{\implies} Mostrar que los isomorfismos preservan clases de conjugación. La clase no trivial más pequeña nos da la biyección $X \to Y$. %% Demostrar!!!!!
  \end{itemize}
\end{proof}

\begin{notation}
  \begin{align*}
    S_n &= \Sym(\{1,\dots,n\}) \\
    S_\oo &= \Sym(\Z)
  \end{align*}
  $\Sym_0(\Z)$ es el grupo de permutaciones de $\Z$ de soporte finito.
\end{notation}

\begin{obs}
  $\abs{S_n} = n!$.
\end{obs}

\begin{definition}
  Decimos que $\sigma \in S_n$ es un \wea{ciclo de largo $k$} si existen elementos $a_1,\dots,a_k \in \{1,\dots,n\}$ tales que
  \begin{itemize}
    \item $\sigma(a_i) = a_{i+1}, \quad \forall i \in \{1,\dots,k-1\}$;
    \item $\sigma(a_k) = a_1$;
    \item $\sigma(x) = x, \quad \forall x \in \{1,\dots,n\} \setminus \{a_1,\dots,a_k\}$.
  \end{itemize}
\end{definition}

\begin{notation}
  El ciclo anterior se denota
  \[\sigma = (a_1,a_2,a_3,\dots,a_k).\]
\end{notation}

\begin{obs}
  La notación no es única.
\end{obs}

\begin{example}
  En $S_5$, tenemos que
  \[(1,2)\c(2,3,5) = (5,1,2,3).\]
\end{example}

\begin{obs}
  Si $\tau = (a_1,\dots,a_k)$ y $\sigma = (b_1,\dots,b_l)$ son tales que
  \[\{a_1,\dots,a_k\} \cap \{b_1,\dots,b_l\} = \varnothing.\]
  Entonces $\tau \circ \sigma = \sigma \circ \tau$.
\end{obs}

\begin{obs}
  Toda permutación en $S_n$ se escribe como producto finito de ciclos disjuntos.
\end{obs}

\begin{definition}
  Una \wea{transposición} es un ciclo de largo 2.
\end{definition}

\begin{obs}
  Todo ciclo en $S_n$ se escribe como producto de transposiciones:
  \[(a_1,\dots,a_k) = (a_1,a_k)(a_1,a_{k-1}) \dots (a_1,a_3)(a_1,a_2).\]
\end{obs}

\begin{obs}
  Toda permutación en $S_n$ es producto de transpocisiones.
\end{obs}

\begin{definition}
  Sea $G$ grupo y $S \subset G$. El \wea{subgrupo de $G$ generado por $S$} es
  \[\gen{S} = \{g \in G \colon \exists n \geq 0 \colon g = s_1,\dots,s_n, \quad s_1,\dots,s_n \in S \cup S^{-1}\}.\]
  Diremos que $S \subset G$ es \wea{generador} si $\gen{S} = G$.
  Diremos que $G$ es \wea{finitamente generado} si $\exists S \subset G$ finito y generador.
\end{definition}

\begin{obs}
  Si $G$ es finitamente generado, entonces $G$ es contable (finito o numerable).
\end{obs}

\begin{obs}
  $S_n$ es generado por las transposiciones.
\end{obs}

\begin{exercise}
  Mostrar que $\Sym_0(\Z)$ no es finitamente generado.
\end{exercise}

\begin{theorem}
  Si $\sigma \in S_n$ se puede escribir como producto de una cantidad par de transposiciones, entonces no es posible escribirla como producto de una cantidad impar de transposiciones.
\end{theorem}
\begin{proof}
  Primero demostrar para la identidad. %% demostrar
  Si
  \[\tau = t_1 t_2 \dots t_k = r_1 r_2 \dots r_l\]
  con $t_i, r_i$ transposiciones, entonces
  \[\id = t_1 \dots t_k r_l^{-1} \dots r_1^{-1}
  \implies k \equiv l \pmod 2.\]
\end{proof}

\begin{definition}
  Consideremos la función $\sgn \colon S_n \to \Z$ dada por
  \begin{itemize}
    \item $\sgn(\sigma) = 1$ si $\sigma$ es permutación par;
    \item $\sgn(\sigma) = -1$ si $\sigma$ es permutación impar.
  \end{itemize}
  El \wea{grupo alternante} es
  \[A_n = \ker(\phi) = \{\sigma \in S_n \colon \sgn(\sigma) = 1\}.\]
\end{definition}

\begin{definition}
  El \wea{grupo diedral} $D_n$ es el grupo de simetrias el polígono regular de $n$ lados.
\end{definition}

\begin{example}
  $D_3$ es el grupo de simetrías de un triángulo equilátero. Si denotamos los vértices como 1, 2 y 3, entonces $D_3$ permuta estos números entre si, por lo que $D_3 \leqslant S_3$, y
  \[D_3 = \{\id, (123), (132), (13), (12), (23)\}.\]
  Similarmente, $D_4 \leqslant S_4$, pero no escribiré sus elementos.
\end{example}

\newpage
\subsection{Clases laterales}

\begin{definition}
  Sea $G$ un grupo, $H \leqslant G$. Una \wea{clase lateral} (\wea{izquierda} o \wea{derecha}) es
  \[gH = \{gh \colon h \in H\}, \quad Hg = \{hg \colon h \in H\}\]
  para algún $g \in G$.
\end{definition}

\begin{proposition}
  Las clases laterales de $H \leqslant G$ sorman una partición de $H$:
  \begin{enumerate}[(i)]
    \item $\bigcup_{g\in G} gH = G$;
    \item $gH \cap g'H \neq \varnothing \implies gH = g'H$.
  \end{enumerate}
\end{proposition}
\begin{proof}
  $\bigcup_{g \in G} gH = G$ es trivial.

  Supongamos que $gH \cap g' H \neq \varnothing$, entonces $\exists h,h' \in H$ tales que $gh = g'h'$. Luego $g' = gh(h')^{-1} \in gH$. Luego $\forall h'' \in H$, tenemos que $g'h'' \in gH$, y entonces $g'H \subseteq gH$. Análogo para $gH \subseteq g'H$.
\end{proof}

\begin{obs}
  Supongamos que $H \leqslant G$ y $\abs{H} < \oo$, entonces $\forall g \in G$, $|gH| = |H|$. Si $G$ es finito, entonces
  \[G = \bigcupdot_{g \in L} gH \implies \abs{G} = \abs{L} \c \abs{H}.\]
\end{obs}

\begin{definition}
  Decimos que $L \subset G$ es un \wea{conjunto de representantes izquierdos} (\wea{derechos}) de $H \leqslant G$ si
  \[G = \bigcupdot_{g \in L} gH = \bigcupdot_{g \in L} Hg.\]
\end{definition}

\begin{theorem}
  [Lagrange]
  Sea $G$ un grupo finito. Si $H \leqslant G$, entonces $|H|$ divide a $|G|$.
\end{theorem}

\begin{definition}
  Un subgrupo $G \leqslant G$ se dice \wea{normal} y se denota $H \normal G$ si $\forall h \in H$ y $g \in G$,
  \[ghg^{-1} \in H.\]
\end{definition}

\begin{proposition}
  $H \normal G \iff \forall g \in G, \ gH = Hg$.
\end{proposition}

\begin{obs}
  Si $G$ es abeliano, entonces $H \leqslant G \implies H \normal G$.
\end{obs}

\begin{example}
  Todos los subgrupos de $(\Z,+)$ son de la forma $n\Z$ con $n \geq 1$. Como $\Z$ es abeliano, estos son subgrupos normales. Las clases laterales son de la forma $k + n\Z$ con $k \in \{0,\dots,n-1\}$. Notemos que
  \[(k + n\Z) + (k'+n\Z) = (k+k' \pmod n + n\Z)\]
\end{example}

\fecha{Clase 3. Lunes 16 de Marzo}

\begin{definition}
  Dado $K \normal G$, el \wea{grupo cociente} es
  \[\frac{G}{K} = \{gK \colon g \in G\}\]
  con la operación
  \[(gK) \c (hK) = ghK.\]
\end{definition}

\begin{definition}
  Si $H \leqslant G$, el \wea{índice de $H$ en $G$} es $$[G \colon H] = \abs{\frac{G}{H}}.$$
\end{definition}

\begin{proposition}
  Sea $\phi \colon G \to H$ homomorfismo de grupos. Entonces $\ker(\phi) \normal G$.
\end{proposition}
\begin{proof}
  Sea $g \in \ker(\phi)$ y $h \in G$, entonces
  \[\phi(hgh^{-1}) = \phi(h) \phi(g) \phi(h)^{-1} = \phi(h) \c e_H \c \phi(h)^{-1} = \phi(h) \phi(h)^{-1} = e_H\]
  y entonces $hgh^{-1} \in \ker(\phi)$, para todo $g \in \ker(\phi)$ y todo $h \in G$. Por lo tanto $\ker(\phi) \normal G$.
\end{proof}

\begin{proposition}
  Sea $K \normal G$. Consideremos la proyección canónica $$\phi \colon G \to \frac GK$$ dada por $\phi(g) = gK$. Entonces $\phi$ es un epimorfismo y $\ker(\phi) = K$.
\end{proposition}

\begin{obs}
  Tenemos una correspondencia entre subgrupos normales y núcleos de homomorfismos. %% Escribir mejor
\end{obs}

\begin{theorem}[\wea{Primer Teorema del Isomorfismo}]
  Sean $G,H$ grupos y $\phi \colon G \to H$ homomorfismo. Entonces
  \[\im(\phi) \simeq \frac{G}{\ker(\phi)}.\]
\end{theorem}
\begin{proof}
  Sea $K = \ker(\phi)$. Definimos $$\psi \colon \frac G K \to \im(\phi)$$ dada por $\psi(gK) = \phi(g)$. Veamos que $\psi$ es un isomorfismo.
  \begin{itemize}
    \item \wea{Bien definido:} Sean $g,h \in G$ tal que $gK = hK$. Luego
    $$gh^{-1} \in K = \ker(\phi)$$
    $$\implies \phi(gh^{-1}) = 1_H$$
    $$\implies \psi(gK) = \phi(g) = \phi(gh^{-1}) \phi(h) = \phi(h) = \psi(hK).$$
    Por lo tanto $\psi$ está bien definida.
    \item \wea{Homomorfismo:} Sean $g,h \in G$, entonces
    \[\psi((gK)(hK)) = \psi(ghK) = \phi(gh) = \phi(g) \phi(h) = \psi(gK) \psi(hK).\]
    \item \wea{Inyectivo:} Sea $g \in G$ tal que $gK \in \ker(\psi)$, entonces
    \[1_H = \psi(gK) = \phi(g) \implies g \in \ker(\phi) = K \implies gK = 1_{g/K}.\]
    Por lo tanto $\ker(\psi) = 1$, y entonces $\psi$ es inyectivo.
    \item \wea{Sobreyectivo:} Trivial. %% demostrar
  \end{itemize}
\end{proof}

\begin{corollary}
  Si $\phi \colon G \to H$ es epimorfismo, entonces $$H \simeq \rac{G}{\ker(\phi)}.$$
\end{corollary}

\begin{example}
  Sean $m,n$ enteros positivos. Pruebe que $\exists \phi \colon \zmz n \to \zmz m$ epimorfismo si y solo si $m$ divide a $n$.
\end{example}
\begin{proof}\
  \begin{itemize}
    \item[\boxed{\implies}] Supongamos que existe un epimorfismo $\phi \colon \zmz n \to \zmz m$, entonces
    \[\znz m \simeq \frac{\zmz n}{\ker(\phi)}\]
    \[\implies m
    = \abs{\znz m}
    = \abs{ \frac{\zmz n}{\ker(\phi)} }
    = \frac{n}{\abs{\ker(\phi)}}\]
    y entonces $n \mid m$.

    \item[\boxed{\impliedby}] Ejercicio. %% demostrar
  \end{itemize}
\end{proof}

\begin{obs}
  En notación moderna, el Primer Teorema de Isomorfismo se escribe: Existe la secuencia exacta corta
  \[1 \to N \to G \to H \to 1.\] %% mostrar que esto equivale al 1er teo de isom
\end{obs}

\begin{definition}
  Sean $H,N \leqslant G$, definimos
  \[HN = \{g \in G \colon \exists h \in H, \exists n \in N, g = hn\}.\]
\end{definition}

\begin{proposition}
  $HN \leqslant G \iff H \normal G \ \lor\ N \normal G$.
\end{proposition}
\begin{proof}
  \
  \begin{itemize}
    \item[\boxed{\impliedby}]
    Supongamos que $N \normal G$. Sean $h,h' \in H$ y $n,n' \in N$. Entonces
    \[(h'n') \c (hn)^{-1}
    = h' n' n^{-1} h^{-1}
    = \underbrace{h' h^{-1}}_{\in H} \underbrace{h n' n^{-1} h^{-1}}_{\in N} \in HN.\]
    \item[\boxed{\implies}] Pendiente %% buscar contraejemplo.
  \end{itemize}
\end{proof}

\begin{theorem}
  [\wea{Segundo Teorema del Isomorfismo}]
  Sea $G$ grupo, $N \normal G$ y $H \leqslant G$. Entonces
  \[\frac{H}{H \cap N} \simeq \frac{HN}{N}.\]
\end{theorem}
\begin{proof}
  Primero notemos que $H \cap N \normal H$ y que $N \normal HN$. %% demostrar
  Definimos $$\phi \colon H \to \frac{HN}{N}$$ mediante $$\phi(h) = hN.$$
  \begin{itemize}
    \item \wea{$\phi$ es homomorfismo:} Sean $h,h' \in H$, entonces $$\phi(hh') = hh'N = (hN)(h'N) = \phi(h) \phi(h').$$
    \item \wea{Sobreyectivo:} Trivial. %% demostrar
  \end{itemize}
  Luego $\phi$ es un epimorfismo, y entonces $$\frac{H}{\ker(\phi)} \simeq \frac{HN}{N}.$$ Notemos que
  \begin{align*}
    \ker(\phi) &= \{h \in H \colon \phi(h) = N\} \\
    &= \{h \in H \colon hN = N\} \\
    &= H \cap N.
  \end{align*}
  Entonces $$\frac{H}{H \cap N} = \frac{H}{\ker(\phi)} \simeq \frac{HN}{N}.$$
\end{proof}

\begin{obs}
  En notación moderna: Tenemos las dos secuencias exactas
  \[\begin{tikzcd}
	&& N \\
	1 & {H \cap N} && HN & G \\
	&& H
	\arrow[from=1-3, to=2-4]
	\arrow[from=2-1, to=2-2]
	\arrow[from=2-2, to=1-3]
	\arrow[from=2-2, to=3-3]
	\arrow[from=2-4, to=2-5]
	\arrow[from=3-3, to=2-4]
  \end{tikzcd}\]
\end{obs}

\begin{example}
  Sea $G$ grupo, $H \leqslant G$, $N \normal G$.
  \begin{enumerate}[(1)]
    \item Suponga que $H \cap N = \{1_H\}$ y que $\forall h \in H$, $hnh^{-1} = n$. $$\implies HN \simeq H \times N.$$ En este caso $$H \simeq \frac{H \times N}{\{1_H\} \times N}.$$
    \item Supongamos que $H \normal G, N \normal G, H \cap N = \{1_G\}$. $$\implies HN \simeq H \times N.$$
  \end{enumerate}
\end{example}

\begin{theorem}
  [\wea{Tercer Teorema del Isomorfismo}]
  Sea $G$ grupo y $H,N \normal G$ tales que $N \leqslant H \leqslant G$.
  \[\implies \frac{G}{H} \simeq \frac{G/N}{H/N}.\]
\end{theorem}
\begin{proof}
  Definimos $$\phi \colon \frac GN \to \frac GH$$ dado por $$\phi(gN) = gH.$$ Veamos que es epimorfismo.
  \begin{itemize}
    \item \wea{Bien definido:} Sean $g,g' \in G$ tales que $gN = g'N$. Entonces $g'g^{-1} \in N \subset H$, y luego $g'g^{-1} \in H$ y por lo tanto $$\phi(g'N) = g'H = g'g^{-1}gH = (g'g^{-1}H)(gH) = gH = \phi(gN).$$
    \item \wea{Homomorfismo:} Trivial %% demostrar
    \item \wea{Sobreyectivo:} Trivial %% demostrar
  \end{itemize}
  Por Primer Teorema del Isomorfismo, basta verificar que $$\ker(\phi) = \frac HN.$$
  \begin{align*}
    \ker(\phi) &= \{gN \colon \phi(gN) = 1_{G/H} = H\} \\
    &= \{gN 'colon gH = H\} \\
    &= \{gN \colon g \in H\} \\
    &= \frac HN.
  \end{align*}
\end{proof}

\begin{definition}
  Un grupo $G$ se dice \wea{simple} si los únicos subgrupos normales son $\{1_G\}$ y $G$.
\end{definition}

\begin{exercise}
  Pruebe que si $p$ es primo, entonces $\zmz p$ es simple.
\end{exercise}

\newpage
\subsection{Acciones de Grupo}

\begin{definition}
  Sea $X$ un conjunto y $G$ un grupo. Una \wea{acción (izquierda) de $G$ en $X$} es una función $$\alpha \colon G \times X \to X$$ tal que
  \begin{enumerate}[(1)]
    \item $\alpha(1_G,x) = x, \quad \forall x \in X$;
    \item $\alpha(gh,x) = \alpha(g, \alpha(h,x)), \quad \forall g,h\in G, \forall x \in X$.
  \end{enumerate}
\end{definition}

\begin{notation}
  Una acción $\alpha$ de $G$ en $X$ se denota por $$G \actsby{\alpha} X$$ y $$\alpha(g,x) = g \c x = g \c_\alpha x.$$
\end{notation}

\begin{example}
  $\GL_n(\Z) \actsby\alpha \Z^n$ mediante $\alpha(A,x) = Ax$.
\end{example}

\begin{example}
  $\GL_n(\Z) \acts \R^n/\Z^n$ por multiplicación.
\end{example}

\begin{example}
  $\Sym(X) \actsby\alpha X$ por $\alpha(\sigma,x) = \sigma(x)$.
\end{example}

\begin{proposition}
  Sea $G$ grupo, entonces $\Aut(G) \acts G$ mediante $\phi \c g = \phi(g)$.
\end{proposition}

\begin{example}
  $G \actsby\alpha G$ por conjugación: $$g \c_\alpha h = ghg^{-1}.$$
\end{example}

\begin{example}
  $G \actsby\alpha G$ por traslación izquierda: $$g \c_\alpha h = gh.$$
\end{example}

\begin{obs}
  Hay una correspondencia
  \[
  \set{\text{Acciones $G \acts X$}} \leftrightarrow \Hom(G,\Sym(X))
  \]
  Dada una acción $G \actsby\alpha X$, definimos $\phi_\alpha(g) = \alpha(g,\c) \in \Hom(G,\Sym(X))$.
\end{obs}

\begin{definition}
  Sea $G \acts X$ una acción de grupo.
  \begin{itemize}
    \item La \wea{órbita de $x \in X$} es $$\orb_G(x) = \{g \c x \colon g \in G\}.$$
    \item El \wea{estabilizador de $x \in X$} es $$\stab_G(x) = \{g \in G \colon g \c x = x\}.$$
  \end{itemize}
\end{definition}

\begin{obs}
  $\stab_G(x) \leqslant G$. No es necesariamente un subgrupo normal. %% buscar contraejemplo
\end{obs}

\begin{proposition}
  Dada una acción de grupo $G \acts X$, para todo $x \in X$ tenemos que
  $$\bigcap_{x \in X} \stab_G(x) \normal G.$$
\end{proposition}

\fecha{Clase 4. Miércoles 18 de Marzo}

\begin{definition}
  Decimos que una acción $G \acts X$ es \wea{fiel} si $$\forall g \in G \setminus \{1_G\}, \exists x \in X \colon g\c x \neq x.$$
\end{definition}

\begin{proposition}
  $G \acts X$ es fiel $\ds \iff \bigcap_{x \in X} \stab_G(x) \simeq 1$.
\end{proposition}

\begin{proposition}
  $G \actsby\alpha X$ es fiel $\iff \ker(\phi_\alpha) \simeq 1$.
\end{proposition}

\begin{example}
  Considere la acción $\Z^2 \acts \R^2$ dada por $$(n,m) \c (x_1,x_2) = (x_1 + n, x_2 + n).$$
  Esta acción no es fiel, pues $(0,m) \c x = x$ para todo $n \in \Z$ y todo $x \in \R^2$. Notemos también que
  \[\stab_{\Z^2}(\R^2) = \bigcap_{x \in \R^2} \stab_{\Z^2}(x) = \{0\} \times \Z \leqslant \Z^2.\]
  Notemos que tenemos una acción
  \[\frac{\Z^2}{\{0\} \times \Z} \simeq \Z \acts \R^2, \quad n \c (x_1,x_2) = (x_1 + n, x_2 + n).\]
\end{example}

\begin{definition}
  Decimos que $G \actsby\alpha X$ es \wea{libre} si
  \[\forall g \in G, \forall x \in X \cor{g\c x=x \implies g = 1_G}.\]
\end{definition}

\begin{proposition}
  Son equivalentes:
  \begin{enumerate}[(i)]
    \item $G \actsby\alpha X$ es libre;
    \item $\forall g \in G, \forall x \in X \cor{g \neq 1_G \implies gx \neq  x}$;
    \item $\forall x \in X, \stab_G(x) \simeq 1$;
    \item $\forall g \in G \setminus \{1_G\}$, $\phi_\alpha(g)$ no tiene puntos fijos.
  \end{enumerate}
\end{proposition}

\begin{example}
  La acción izquierda por multiplicación de un grupo en sí mismo es libre.
\end{example}

\begin{example}
  Considere $S^1 = \{z \in \C \colon \abs z = 1\}, \ \alpha \in [0,1)$. Tenemos una acción $\Z \acts S^1$ dada por $$n \c z = z e^{2\pi in\alpha}.$$
  Si $\alpha \in \Q, \alpha = p/q$ con $p,q \geq 1$ enteros y $p < q$, entonces la acción \wea{no} es fiel por
  \[q \c z = ze^{2\pi i \frac{p}{q} q} = ze^{1\pi i p} = z.\]
  Si $\alpha \notin \Q$ y $z \in S^1, \ n \in \Z$ son tales que $n \c z \implies e^{2\pi i n \alpha} = 1$. Por lo tanto $n\alpha \in \Z \implies n=0$. Entonces la acción es libre.
\end{example}

\begin{example}
  Si $X$ es un conjunto con $\abs X > 2$, entonces $\Sym(X) \acts X$ es fiel, pero no es libre.
\end{example}

\begin{definition}
  $G \acts X$ se dice \wea{transitiva} si $\forall x,y \in X, \exists g \in G \colon g\c x = y$.
\end{definition}

\begin{proposition}
  Dada una acción $G \acts X$, son equivalentes:
  \begin{enumerate}[(i)]
    \item $G \acts X$ es transitiva;
    \item $\forall x \in X, \orb_G(x) = X$;
    \item $\exists x \in X \colon \orb_G(x) = X$.
  \end{enumerate}
\end{proposition}

\begin{definition}
  Decimos que un grupo $H$ \wea{se incrusta en un grupo $G$} si $\exists \phi \colon H \to G$ monomorfismo. Se denota $H \embed G$.
\end{definition}

\begin{theorem}
  [Cayley]
  Sea $G$ grupo, entonces $G \embed \Sym(G)$. Si $\abs G = n < \oo$, entonces $G \embed S_n$.
\end{theorem}
\begin{proof}
  Consideremos $\phi \colon G \to \Sym(G)$ dada por $$(\phi(g))\c(h) = gh, \quad \forall g,h \in G.$$
  Este es un monomorfismo.
\end{proof}

\begin{exercise}
  Sea $G$ grupo finito con $\abs G = n$. Pruebe que $$G \embed \GL_n(\Z).$$
\end{exercise}
%% demostrar. primero G \embed \Sym(G) \simeq S_n \embed \GL_n(\Z), este último por matrices de permutación.

\begin{theorem}
  [Órbita--Estabilizador]
  Sea $G \acts X$ una acción y $x \in X$. Entonces existe una biyección
  \[f \colon \orb_G(x) \to \frac{G}{\stab_G(x)}.\]
  En particular, $\abs G < \oo \implies \abs{\orb_G(x)} \c \abs{\stab_G(x)} = \abs{G}$.
\end{theorem}
\begin{proof}
  Definimos $f \colon \orb_G(x) \to G/\stab_G(x)$ por $f(g \c x) = g \stab_G(x)$.
  \begin{itemize}
    \item \wea{Bien definido: }
    Sean $g,h \in G$ tales que $g \c x = h \c x$. Entonces $h^{-1} g \c x = x$, y entonces $h^{-1}g \in \stab_G(x)$. Luego $$f(h \c x) = h \stab_G(x) = g \stab_G(x) = f(g \c x).$$

    \item \wea{Sobreyectivo: } \checkmark %% ????
    \item \wea{ Inyectivo: }
    Sean $g\c x, h \c x \in \orb_G(x)$ tales que $f(g \c x) = f(h \c x)$, entonces $g \stab_G(x) = h \stab_G(x)$. Luego $h^{-1} g \in \stab_G(x)$ y entonces $g\c x = h \c x$.
  \end{itemize}
\end{proof} %% arreglar el como se ve el g\stab

\begin{obs}
  $\stab_G(x)$ no es necesariamente un subgrupo normal de $G$, pero el teorema trabaja solo con las clases laterales, sin estructura de grupo.
\end{obs}

\begin{theorem}
  [Cauchy]
  Sea $G$ grupo finito y sea $p$ primo tal que $p$ divide a $\abs G$. Entonces $\zmz p \embed G$.
\end{theorem}
\begin{proof}
  Consideremos $$X = \{(x_1,\dots,x_p) \in G^p \colon x_1 x_2 \cdots x_p = 1_G\}.$$
  Notemos que $$X = \{(x_1,\dots,x_{p-1}, x_{p-1}^{-1}\cdots x_2^{-1}x_1^{-1}) \colon x_1,\dots,x_{p-1} \in G\}$$ y entonces $\abs X = \abs G^{p-1}$.
  Consideremos la acción $\zmz p \acts X$ por permutación cíclica mediante $$1 \c (x_1,\dots,x_p) = (x_2,x_3, \dots, x_p, x_1).$$
  Notemos que si $(x_1,\dots,x_p) \in X$, entonces
  \begin{align*}
    x_1 x_2 \cdots x_p &= 1_G \\
    \implies x_2 x_3 \cdots x_p x_1 &= x_1^{-1} x_1 x_2 x_2 \cdots x_p x_1 \\
    &= x_1 1_G x_1 = 1_G.
  \end{align*}
  Luego
  \[\abs{X} = \abs{G}^{p-1} = \sum_{x_i \in X/(\Z/p\Z)} |\orb_{\Z/p\Z} (x_i)|.\] %% no entiendo este error, lo arreglaré
  Supongamos por contradicción que $\forall g \in G \setminus \{1_G\}, \quad g^p \neq 1_G$.
  Luego $\forall x = (x_1,\dots,x_p) \neq (1_G,\dots,1_G)$ tenemos que $\abs{\orb_{\zmz p}(x)} = p$. Luego
  \[\abs G^{p-1} = 1 + p \br{\abs{\frac{X}{\zmz p}} -1}.\]
  Tomando $\mod p$ tenemos que $0=1 \mod p$. $\to\gets$
\end{proof}

\begin{theorem}
  Sea $G$ un grupo finito y $p$ primo tal que $\abs G = p^2$. Entonces $G$ es abeliano.
\end{theorem}

\begin{definition}
  Sea $G$ un grupo. La \wea{clase de conjugación de $h \in G$} es
  \[C(h) \coloneqq \{ghg^{-1} \colon g \in G.\]
\end{definition}

\begin{obs}
  Tenemos la acción $G \acts G$ por congugación $g \star h = ghg^{-1}$. Luego $C(h) = \orb_G(h)$ para todo $h \in G$.
\end{obs}

\begin{obs}
  Si $G$ es finito, entonces $$\abs G = \sum_{i \in I} \abs{C(g_i)}$$ donde $\{g_i\}_{i \in I}$ son representantes de las clases de conjugación. Los números $\{\abs{C(g_i)}\}$ se denomina \wea{ecuación de clase}.
\end{obs}

\begin{example}
  Si $G$ es abeliano, entonces la ecuación de clase de $G$ es $$\abs G = \underbrace{1+1+\dots+1}_{\abs G \text{ veces}}.$$
\end{example}

\begin{example}
  $|S_3| = |C(\id)| + |C((12))| + |C((123))| = 1+3+2$.
\end{example}

\begin{obs}
  La ecuación de clase no es invariante completo de conjugación.
\end{obs}

\begin{definition}
  Sea $G$ un grupo. Su \wea{centro} es
  \[Z(G) = \{g \in G \colon gh=hg, \quad \ h \in G\}.\]
\end{definition}

\begin{obs}
  Sea $\phi \colon G \to \Sym(G)$ por conjugación, entonces $Z(G) = \ker(\phi)$. %% no entendí la función
\end{obs}

\begin{definition}
  Sea $G$ grupo y $g \in G$. El \wea{centralizador de $g$} es
  \[\Cent(g) = \{h \in G \colon hg = gh\}.\]
\end{definition}

\begin{obs}
  $\ds Z(G) = \bigcap_{g \in G} \Cent(g)$.
\end{obs}

\newpage %% error de mdframed, quitar eventualmente este salto de pagina
\begin{lemma}
  Sea $G$ grupo finito, $n\geq 1$ entero y $p$ primo tal que $\abs G = p^n$. Entonces $\abs{Z(G)} > 1$.
\end{lemma}
\begin{proof}
  Supongamos que $\abs{Z(G)}=1$, entonces existe una única clase de conjugación con un elemento. Luego toda clase de conjugación $C(g)$ satisface que $\abs{C(g)}$ es divisible por $p$. Luego
  \[p^n = \abs{G} = \sum_{g_i \text{ rep}} \abs{C(g_i)} = 1 +pk, \quad k \in \Z.\]
  Contradicción, pues $0 \neq 1 \pmod p$.
\end{proof}

\fecha{Clase 5. Lunes 23 de Marzo}

\begin{lemma}
  Todo grupo de orden $p^n$ con $n \geq 1$ y $p$ primo tiene centro no trivial.
\end{lemma}

\begin{theorem}
  Todo grupo de orden $p^2$ es abeliano con $p$ primo.
\end{theorem}
\begin{proof}
  Sea $G$ un grupo de orden $p^2$. Tenemos 3 opciones:
  \begin{enumerate}[(i)]
    \item $Z(G) \simeq 1$;
    \item $\abs{Z(G)} = p$;
    \item $Z(G) = G$.
  \end{enumerate}
  La opción (i) no funciona por el lema. Supongamos que $\abs{Z(G)} = p$.
\end{proof}

\newpage
\subsection{Subgrupos Simples}

\begin{definition}
  Un grupo $G$ se dice \wea{simple} si sus únicos subgrupos normales son $\{1_G\}$ y $G$.
\end{definition}

\begin{proposition}
  $G$ es simple si todo cociente de $G$ es isomorfo a $G$ o 1.
\end{proposition}

\begin{example}
  $\zmz p$ es simple para $p$ primo.
\end{example}

\begin{obs}
  Recordemos que dado $n \geq 2$, las permutaciones $\sigma \in S_n$ tienen signo:
  \begin{itemize}
    \item $\sgn(\sigma) = 1$ si $\sigma$ es par;
    \item $\sgn(\sigma) = -1$ si $\sigma$ es impar.
  \end{itemize}
  Esto nos da un homomorfismo de grupos $\phi \colon S_n \to \Z_2$, y entonces $A_n = \ker(\phi)$. Por lo tanto $A_n \normal S_n$.

  Sea $\sigma \in S_n$ un producti de ciclos
  \[\sigma = (a_1,\dots,a_{k_1})(b_1,\dots,b_{k_2})\dots(t_1,\dots,t_{k_l}).\]
  Sea $\tau \in S_n$, entonces
  \[\tau\sigma\tau^{-1} = (\tau(a_1,\dots,a_{k_1})\tau^{-1})(\tau(b_1,\dots,b_{k_2})\tau^{-1}) \dots (\tau(t_1,\dots,t_{k_l})\tau^{-1}).\]
  Luego basta estudiar la conjugación de un ciclo. Sea $\sigma = (a_1 a_2 \dots a_k)$ y $\tau \in S_n$ con $k \leq n$, entonces
  \[\tau\sigma\tau^{-1} = (\tau(a_1), \dots, \tau(a_k)).\]
\end{obs}

\begin{obs}
  La conjugación en $S_n$ preserva la estructura de ciclos.
\end{obs}

\begin{example}
  En $S_{11}$ (escrito en hexadecimal) tomemos
  \[\sigma = (135)(24)(6789)(AB), \quad \tau = (123456789AB)\]
  \[\implies \tau\sigma\tau^{-1} = (246)(35)(789A)(B1).\]
\end{example}

\begin{example}
  Calculemos las clases de conjugación de $S_5$. Primero notemos que $\abs{S_5} = 120$. La estructura de las permutaciones son:
  \begin{enumerate}[(1)]
    \item $\id$ (1 elemento);
    \item $(ab)$ (10 permutaciones);
    \item $(abc)$ (20 permutaciones);
    \item $(abcd)$ (30 permutaciones);
    \item $(abcde)$ (24 permutaciones);
    \item $(ab)(cd)$ (15 permutaciones);
    \item $(ab)(cde)$ (20 permutaciones).
  \end{enumerate}
  Como la conjugación preserva la estructura de los ciclos, tenemos que las clases de conjugaciones coinciden con estos números, y entonces la ecuación de clase es
  \[120 = 1 + 10 + 15 + 20 + 20 + 24 + 30.\]
  Vamos a estudiar las clases de conjugación en $A_5$. Sus elementos son de la forma
  \[\id, \quad (abc) = (ac)(ab), \quad (abcde) = (ae)(ad)(ac)(ab), \quad (ab)(cd).\]
  Veamos las clases de conjugación de este tipo de elementos.
  \begin{itemize}
    \item Los elementos de la forma $(abc)$ forman una única clase de conjugación:
    Supongamos que $(abc)$ y $(a'b'c')$ son conjugados por $\tau \in S_n$ con $n \geq 5$. Sean $e,f \in \{1,\dots,n\} \setminus \{a,b,c\}$ con $e \neq f$. Definimos $\tau' = \tau(ef)$. Luego
    \[\tau'(abc)\tau'^{-1} = (\tau'(a), \tau'(b), \tau'(c))
    = (\tau(a),\tau(b),\tau(c)) = \tau(abc)\tau^{-1}.\]
    Luego $\abs{C_{A_n}((123))} = 20$.

    \item Para $(ab)(cd)$, consideremos el siguiente lema:
    \begin{lemma}
      Sea $G$ grupo finito y $G \acts X$ transitiva. Sea $N \normal X$ con $[G \colon N] = 2$ tal que $N \acts X$ no es transitiva.
      Entonces $N \acts X$ tiene 2 órbitas de misma cardinalidad.
    \end{lemma}
    Aplicamos con $G = S_5$, $X = C_{S_5}((12)(34))$ con $N = A_5$ y $G \acts X$ por conjugación.
    Luego si $C_{A_5}((12)(34)) \neq C_{S_5}((12)(34))$, entonces $\abs{C_{A_5}((12)(34))} = 15/2$, lo cual no puede ser. Entonces $\abs{C_{A_5}((12)(34))} = 15$.

    \item Para los 5--ciclos, veremos que $(12345)$ y $(21345)$ no son conjugados. Supongamos que sí lo son, entonces existe $\tau \in S_5$ tal que $(21345) = \tau(12345)\tau^{-1} = (\tau(1), \dots, \tau(5))$.
    Entonces $\tau = (12) \sigma$ con $\sigma \in \gen{(12345)}$. Por lo tanto $\tau \in (12)A_5$. Luego
    \[C_{A_5}((12345)) \cap C_{A_5}((21345)) = \varnothing.\]
    Por el lema, ésta última clase de conjugación se separa en 2.
  \end{itemize}
  Luego, la ecuación de clase de $A_5$ es $$\abs{A_5} = 1 + 12 + 12 + 15 + 20.$$
\end{example}

\fecha{Clase 6. Miércoles 25 de Marzo}

\begin{theorem}
  $A_5$ es simple.
\end{theorem}
\begin{proof}
  Sea $N \normal A_5$.
  Primero, como $N$ es normal, toda clase de conjugación $C$ cumple $C \subset N$ o $C \cap N = \varnothing$.

  También, como $N \leqslant A_5$, tenemos que $\abs N$ divide a 60. Luego
  \[\abs{N} \in \{1,2,3,4,5,6,10,12,15,20,30,60\}.\]
  Luego $$\abs N = 1 + 12a_1 + 12a_2 + 15a_3 + 20a_4$$ con $a_i \in \{0,1\}$. Por lo tanto $\abs N \in \{1,60\}$.
\end{proof}

\begin{exercise}
  Muestre que $A_4$ no es simple.
\end{exercise}

\begin{theorem}
  $A_n$ es simple para $n \geq 5$.
\end{theorem}
\begin{proof}
  Por inducción, supongamos que $A_{n-1}$ es simple para algún $n \geq 6$.
  Tomamos $N \normal A_n$ con $N \neq A_n$, veamos que $N \simeq 1$.

  Notemos que con la acción $A_n \acts \{1,\dots,n\}$, tenemos que $$\stab_{A_n}(i) \simeq A_{n-1}, \quad \forall i \in \{1,\dots,n\}.$$
  Luego $$N \cap \stab_{A_n}(i) \normal \stab_{A_n}(i) \simeq A_{n-1}.$$
  Como $A_{n-1}$ es simple (por hipótesis), tenemos que
  \[N \cap \stab_{A_n}(i) \in \set{\set{\id}, \stab_{A_{n-1}}(i)}.\]
  Si $\stab_{A_n}(i) \subset N$, entonces $\forall \sigma \in A_n$,
  \[\stab_{A_n}(\sigma(i)) = \sigma \stab_{A_n}(i) \sigma^{-1} \subset \sigma N \sigma^{-1} = N\]
  \[\implies \bigcup_{1 \leq i \leq n} \stab_{A_n}(i) \subset N.\]
  En particular, para $n \geq 4$, $N$ contiene todos los 3--ciclos.
  \begin{lemma}
    Para $n \geq 3$, $A_n$ está generado por los 3--ciclos.
  \end{lemma}
  Luego $A_n \subset N$. Por lo tanto $N = A_n$, contradicción. Deducimos entonces que
  \[N \cap \stab_{A_n}(i) \simeq 1, \quad \forall i.\]
  Por lo tanto $\forall \sigma \in N \setminus \{\id\}$, $\sigma$ no tiene puntos fijos, es decir: $N \acts \{1,\dots,n\}$ es libre.

  Sea $\sigma \in N \setminus \{\id\}$. Veamos que existe $\tau \in A_n$ (con $n \geq 6$) tal que
  \[\tau(1) = 1 \quad\land\quad \tau(\sigma(1)) = \sigma(1) \quad\land\quad \tau\sigma \neq \sigma\tau.\]
\end{proof}

\fecha{Clase 7. Lunes 30 de Marzo}

El objetivo ahora es dar un ejemplo de un grupo simple infinito.

Recordemos que
\[S_\oo = \Sym_0(\Z) = \{\sigma \in \Sym(\Z) \colon \abs{\sop(\sigma)} < \oo\}\]
donde $$\sop(\sigma) = \{n \in \Z \colon \sigma(n) \neq n\}.$$

Notemos que
\[\Sym_0(\Z) = \bigcup_{n \geq 1} \Sym(\{-n, \dots, n\}).\]
Luego $\sgn \colon \Sym_0(\Z) \to \Z_2$ está bien definida y $A_\oo = \ker(\sgn)$ está bien definido.

\begin{obs}
  $A_\oo$ está generado por los 3--ciclos.
\end{obs}

\begin{exercise}
  Pruebe que $A_\oo$ no es finitamente generado.
\end{exercise}

\begin{definition}
  Sea $G$ un grupo y sean $x,y \in G$. El \wea{conmutador de $x$ e $y$} es
  \[[x,y] = xyx^{-1}y^{-1}.\]
  El \wea{subgrupo derivado} (o \wea{subgrupo conmutador}) de $G$ es
  \[G' = [G,G] = \gen{[x,y] \colon x,y \in G}.\]
\end{definition}

\begin{definition}
  Sea $G$ grupo y $H \leqslant G$. Decimos que $H$ es \wea{característico en $G$} si $$\forall \vphi \in \Aut(G), \quad \vphi(H) = H.$$
\end{definition}

\begin{obs}
  Si $H \leqslant G$ es característico, entonces es normal en $G$. Esto porque $\Inn(G) \leqslant \Aut(G)$, y tenemos que $gHg^{-1} = H$.
\end{obs}

\begin{proposition}
  $G'$ es característico en $G$.
\end{proposition}
\begin{proof}
  Sea $\vphi \in \Aut(G)$ y $g \in G'$, entonces $$g = c_1 c_2 \dots c_n$$ con $c_i = [x_i,y_i]$ y $x_i,y_i \in G$. Luego
  \[\vphi(g) = \vphi(c_1) \dots \vphi(c_n)\]
  y notemos que
  \begin{align*}
    \vphi(c_i) &= \vphi([x_i,y_i]) \\
    &= \vphi(x_i y_i x_i^{-1} y_i^{-1}) \\
    &= \vphi(x_i) \vphi(y_i) \vphi(x_i)^{-1} \vphi(y_i)^{-1} \\
    &= [\vphi(x_i), \vphi(y_i)] \in G'
  \end{align*}
  y entonces $\vphi(g) \in G'$. Por lo tanto $\vphi(G') \subseteq G'$
\end{proof}

\begin{exercise}
  Sea $G$ grupo, $N \normal G$ y sea $H \leqslant N$ característico en $N$. Entonces $H \normal G$.
\end{exercise}

\begin{definition}
  Sea $G$ un grupo. Su \wea{abelianización} es el cuociente
  \[G^\ab = \frac{G}{G'}.\]
\end{definition}

\begin{obs}
  $G^\ab$ es abeliano, y si $N \normal G$ tal que $G/N$ es abeliano, entonces $G' \subset N$.
  \[\begin{tikzcd}
    G & \\
    & {G^\ab} \\
    Q
    \arrow[from=1-1, to=2-2]
    \arrow[from=1-1, to=3-1]
    \arrow[from=2-2, to=3-1]
  \end{tikzcd}\]
\end{obs}

\begin{proposition}
  $$A_\oo = [S_\oo,S_\oo] = S_\oo',$$
  y en particular
  $$A_n = [S_n, S_n] = S_n'.$$
\end{proposition}
\begin{proof}
  Sea $g \in S_\oo'$, entonces $g = c_1 \dots c_n$ con $c_i = [\sigma_i, \tau_i]$ con $\sigma_i, \tau_i \in S_\oo$.
  \[\sgn(g) = \prod_{i=1}^n \sgn(c_i) = 1\]
  \[\sgn(c_i) = \sgn(\sigma_i) \sgn(\tau_i) \sgn(\sigma_i)^{-1} \sgn(\tau_i)^{-1} = 1\]
  y entonces $S_\oo' \subset A_\oo$.
  Para ver  que $A_\oo \subset S_\oo'$, basta ver que todo 3--ciclo es un conmutador:
  \[(abc) = (ac)(bc)(ac)(bc) = [(ac),(bc)].\]
\end{proof}

\begin{lemma}
  Sea $N \normal S_\oo$, $N \neq \{\id\}$. Entonces $A_\oo \subset N$.
\end{lemma}
\begin{proof}
  Sea $\sigma \in N \setminus \{\id\}$.
  \begin{claim}
  Existen $n_1,n_2$ suficientemente grandes tales que
  \[(-n_1, \dots, n_2) \in N.\]
  \end{claim}
  \begin{proof}
    Sea $\sigma = \sigma_1 \dots \sigma_k$ su descomposición en ciclos disjuntos. Escojo $\tau_1 \in S_\oo$ tal que
    \[\sop(\tau_1 \sigma \tau_1^{-1}) \subset - \N, \quad \forall i < k. \tag{$*$}\]
    Tenemos que $\tau_1 \sigma_k \tau_1^{-1} = (1,2,3, \dots, j)$ con $j = \abs{\sigma_k}$.
    Escojo $\tau_2$ tal que
    \[\sop(\tau_2 \sigma_i \tau_2^{-1}) \subset -\N. \tag{$*$}\]
    $\tau_2 \sigma_k \tau_2^{-1} = (j,\dots, 2j-1)$.

    Repitieno muchas veces, obtenemos
    \[\tau_1 \sigma \tau_1^{-1} \dots \tau_n \sigma \tau_n^{-1} = \beta(j,\dots, Mj - (M - 1))\]
    con $\sop(\beta) \subset -\N$. Reconjugando obtengo $\tau$.
  \end{proof}
  Sean $a,b \in S_\oo$. Sea $k \geq 1$ tal que
  \[\sop(a) \cup \sop(b) \subset \{-k,\dots,k\}.\]
  Escojo $\tau \in N$ que permuta cíclicamente un intervalo que contiene a $$\{-k,\dots,3k+1\}.$$
  Sea $l = 2k+1$, entonces
  \[\sop(\tau^l a \tau^{-l}) \cap \sop(b) = \varnothing.\]
  Por lo tanto conmutan.

  Sea $g = [[a,\tau^l],b]$.
  \begin{obs}
    $g \in N$.
  \end{obs}
  \begin{exercise}
    Si $N \normal G$, $h \in N$, $g \in G$, entonces $[g,h] \in N$.
  \end{exercise}
  \begin{obs}
    \begin{align*}
      g &= [[a,\tau^l], b] \\
      &= a \tau^l a^{-1} \tau^{-l} b \tau^l a \tau^{-l} a^{-1} b^{-1} \\
      &= aba^{-1}b^{-1} = [a,b]
    \end{align*}
    y entonces $[a,b] \in N$.
  \end{obs}
  Por lo tanto $A_\oo = A_\oo' \subset N$.
\end{proof}

\begin{obs}
  Si $\{\id\} \neq N \normal S_\oo' = A_\oo$, entonces $[S_\oo',S_\oo'] \subset N$.
\end{obs}

\newpage %% error de mdframed, posiblemente quitarlo eventualmente
\begin{obs}
  $[S_\oo',S_\oo']$ no es trivial.
\end{obs}
\begin{proof}
  Notemos que
  \[[(123),(234)] = (123)(234)(321)(432) = (14)(23) \neq \id.\]
\end{proof}

\begin{theorem}
  $A_\oo$ es simple.
\end{theorem}
\begin{proof}
  Sea $G = S_\oo$. Como $G^{(2)} = (G')'$ es característico en $A_\oo = G'$ y $G' \normal G = S_\oo$,
  \[\implies \{\id\} \neq G^{(2)} = [S_\oo',S_\oo'] \normal S_\oo.\]
  Luego $A_\oo \subset [A_\oo,A_\oo]$, y entonces $A_\oo = [A_\oo,A_\oo]$.

  Sea $\{\id\} \neq N \normal A_\oo$. Por el lema, tenemos que $A_\oo \subset N$, y entonces $N = A_\oo$.
\end{proof}

\newpage
\subsection{Introducción a la Teoría Geométrica de Grupos}

Vamos a ver solo el caso de grupos contables.

\begin{definition}
  Un \wea{grafo dirigido} es un par $(V,E)$ donde $V \neq \varnothing$ es un conjunto (\wea{vértices}) y $E \subset V \times V$ (\wea{aristas}).
\end{definition}

\begin{example}
  $V = \{1,2,3\}, E = \{(1,1), (1,2), (2,3), (3,2), (3,1), (3,3)\}$. El grafo $(V,E)$ se representa como:
  \[\begin{tikzcd}
    1 && 2 \\
    \\
    & 3
    \arrow[from=1-1, to=1-1, loop, in=100, out=170, distance=10mm]
    \arrow[from=1-1, to=1-3]
    \arrow[shift left=2, curve={height=-6pt}, from=1-3, to=3-2]
    \arrow[from=3-2, to=1-1]
    \arrow[shift right, curve={height=-6pt}, from=3-2, to=1-3]
    \arrow[from=3-2, to=3-2, loop, in=305, out=235, distance=10mm]
  \end{tikzcd}\]
\end{example}

\begin{definition}
  Dado un grafo $(V,E)$, un \wea{camino} es una secuencia $v_1,\dots,v_n$ de vértices tal que $(v_i,v_{i+1}) \in E$, $\forall i$.
\end{definition}

\begin{definition}
  Un grafo $(V,E)$ es \wea{fuertemente conexo} si $\forall u,v \in V$ existe un camino de $u$ a $v$.
\end{definition}

\begin{definition}
  Sea $G = (V,E)$ un grafo. Si $\forall e = (a,b) \in E$, existe $\bar e = (b,a) \in E$, y además $G$ es conexo, entonces la \wea{distancia entre $u$ y $v$} es el largo del camino más corto entre $u,v$. Se denota $d(u,v)$.
\end{definition}

\begin{obs}
  La distancia en un grafo conexo es una métrica.
\end{obs}

\begin{definition}
  Sea $G$ un grupo y $S \subset G$. El \wea{grafo de Cayley de $G$ con respecto a $S$} es el grafo
  \[\Cay(G,S) = (G,\{(g,gs) \colon s \in S , g \in G\}).\]
\end{definition}

\begin{example}
  $G = \Z$, $S=\{1\}$, tenemos $E = \{(n,n+1) \colon n \in \Z\}$.
  \[\begin{tikzcd}
    \cdots & \bullet & \bullet & \bullet & \bullet & \cdots
    \arrow[from=1-1, to=1-2]
    \arrow[from=1-2, to=1-3]
    \arrow[from=1-3, to=1-4]
    \arrow[from=1-4, to=1-5]
    \arrow[from=1-5, to=1-6]
  \end{tikzcd}\]
\end{example}

\begin{example}
  $G = \Z^2$, $S = \{(0,0), (1,0)\}$.
  \[\begin{tikzcd}
    \cdots & \bullet & \bullet & \bullet & \bullet & \cdots \\
    \cdots & \bullet & \bullet & \bullet & \bullet & \cdots \\
    \cdots & \bullet & \bullet & \bullet & \bullet & \cdots
    \arrow[from=1-1, to=1-2]
    \arrow[from=1-2, to=1-2, loop, in=55, out=125, distance=10mm]
    \arrow[from=1-2, to=1-3]
    \arrow[from=1-3, to=1-3, loop, in=55, out=125, distance=10mm]
    \arrow[from=1-3, to=1-4]
    \arrow[from=1-4, to=1-4, loop, in=55, out=125, distance=10mm]
    \arrow[from=1-4, to=1-5]
    \arrow[from=1-5, to=1-5, loop, in=55, out=125, distance=10mm]
    \arrow[from=1-5, to=1-6]
    \arrow[from=2-1, to=2-2]
    \arrow[from=2-2, to=2-2, loop, in=55, out=125, distance=10mm]
    \arrow[from=2-2, to=2-3]
    \arrow[from=2-3, to=2-3, loop, in=55, out=125, distance=10mm]
    \arrow[from=2-3, to=2-4]
    \arrow[from=2-4, to=2-4, loop, in=55, out=125, distance=10mm]
    \arrow[from=2-4, to=2-5]
    \arrow[from=2-5, to=2-5, loop, in=55, out=125, distance=10mm]
    \arrow[from=2-5, to=2-6]
    \arrow[from=3-1, to=3-2]
    \arrow[from=3-2, to=3-2, loop, in=55, out=125, distance=10mm]
    \arrow[from=3-2, to=3-3]
    \arrow[from=3-3, to=3-3, loop, in=55, out=125, distance=10mm]
    \arrow[from=3-3, to=3-4]
    \arrow[from=3-4, to=3-4, loop, in=55, out=125, distance=10mm]
    \arrow[from=3-4, to=3-5]
    \arrow[from=3-5, to=3-5, loop, in=55, out=125, distance=10mm]
    \arrow[from=3-5, to=3-6]
  \end{tikzcd}\]
\end{example}

\begin{obs}
  Supongamos dos casos ahora:
  \begin{enumerate}
    \item $S$ es simétrico, i.e.: $S = S^{-1}$;
    \item $G$ es generado por $S$.
  \end{enumerate}
  Entonces $\Cay(G,S)$ es conexo y $E \subset G \times G$ es simétrica. Además, $d_S$ induce una distancia en $G$.

  Además, si $S$ es finito y $1_G \in S$, entonces
  \[\forall g \in G, \forall r \geq 1, \quad \abs{B_S(g,r)} \leq \abs{S}^r\]
  donde $B_S(g,r) = \{h \in G \colon d_S(g,b) \leq r\}$.

  Es decir, $\Cay(G,S)$ es \wea{localmente finito}.
\end{obs}

\begin{definition}
  Sea $G$ grupo y $S \subset G$ generador finito y simétrico. Para $g \in G$ definimos
  \[\abs{g}_S = \min\{n \in \N \colon g = s_1,\dots,s_n, \quad s_i \in S\}.\]
  y definimos $\abs{1_G}_S = 0$.
\end{definition}

\fecha{Clase 8. Miércoles 1 de Abril}

\begin{definition}
  Definimos la \wea{distancia entre $g$ y $h$ sobre $S$} como
  \[d_S(g,h) = \abs{g^{-1} h}_S.\]
\end{definition}

\begin{obs}
  $d_S$ es una distancia y corresponde al largo del camino más corto de $g$ a $h$ en $\Cay(G,S)$.
\end{obs}

\begin{proposition}
  $\forall x, y,g \in G$,
  \[d_S(gx,gy) = d_S(x,y).\]
  i.e.: $d_S$ es invariante bajo traslaciones izquierdas.
\end{proposition}

\begin{obs}
  La elección de generadores no es canónica.
\end{obs}

\begin{example}
  Consideremos $\Z$ con $S_1 = \{-1,0,1\}$. Tenemos que $\Cay(\Z,S_1)$ se ve de la forma:
  \[\begin{tikzcd}
    \cdots & \bullet & \bullet & \bullet & \bullet & \cdots
    \arrow[from=1-1, to=1-2]
    \arrow[from=1-2, to=1-3]
    \arrow[from=1-3, to=1-4]
    \arrow[from=1-4, to=1-5]
    \arrow[from=1-5, to=1-6]
  \end{tikzcd}\]
  Si $S_2 = \{-3,-2,0,2,3\}$, entonces $\Cay(\Z,S_2)$ es
  \[\begin{tikzcd}
    \cdots & {-3} & {-2} & {-1} & 0 & 1 & 2 & 3 & \cdots
    \arrow[curve={height=30pt}, from=1-1, to=1-3]
    \arrow[curve={height=-30pt}, from=1-1, to=1-4]
    \arrow[curve={height=30pt}, from=1-2, to=1-4]
    \arrow[curve={height=-30pt}, from=1-2, to=1-5]
    \arrow[curve={height=30pt}, from=1-3, to=1-5]
    \arrow[curve={height=-30pt}, from=1-3, to=1-6]
    \arrow[curve={height=30pt}, from=1-4, to=1-6]
    \arrow[curve={height=-30pt}, from=1-4, to=1-7]
    \arrow[curve={height=30pt}, from=1-5, to=1-7]
    \arrow[curve={height=-30pt}, from=1-5, to=1-8]
    \arrow[curve={height=30pt}, from=1-6, to=1-8]
    \arrow[curve={height=-30pt}, from=1-6, to=1-9]
    \arrow[curve={height=30pt}, from=1-7, to=1-9]
  \end{tikzcd}\]
\end{example}

\begin{proposition}
  Sea $G$ un grupo finitamente generado y sean $S,S'$ dos generadores finitos y simétricos. Entonces $\exists c>0, \forall g,h \in G$
  \[c^{-1} d_S(g,h) \leq d_{S'}(g,h) \leq c d_S(g,h).\]
\end{proposition}
\begin{proof}
  Notemos que como $S$ es generador, todo elemento de $S'$ se puede escribir como producto de elementos de $S$. Definimos
  \[C_1 = \max_{s' \in S'} \abs{s'}_S.\]
  Luego si $g \in G$,
  \[\abs{g}_{S'} \leq C_1 \abs{g}_S.\]
  En efecto, sea $n = \abs{g}_{S'}$, entonces existen $s_1',\dots,s_n' \in S$ tales que
  \[g = s_1' \c \dots \c s_n'.\]
  Para todo $i \leq n$, existen $s_i^1, \dots, s_i^C \in S$ tales que
  \[s_i' = s_i^1 \cdots s_i^{C_1}.\]
  Luego
  \[g = \prod_{\substack{i \leq n \\ j \leq C_1}} s_i^j.\]
  Luego
  \[\abs{g}_S \leq C_1 \cdot n = \abs{g}_{S'} C_1.\]
  Luego $\forall g,h \in G$,
  \[C_1^{-1} \abs{g^{-1} h}_S \leq \abs{g^{-1} h}_{S'}.\]
  Intercambiando $S$ con $S'$ obtenemos $C_2 > 0$ tal que $\abs{g^{-1}h}_{S'} \leq C_2 \abs{g^{-1} h}_S.$
  Luego $C = \max{C_1,C_2}$ cumple lo pedido.
\end{proof}

\begin{notation}
  Fijo $G$ y $S$. Definimos, para $g \in G$ y $r \in \N$,
  \[B_S(g,r) \coloneqq \{h \in G \colon d_S(g,h) \leq r\}.\]
  \[B_S(r) \coloneqq B_S(1_G,r).\]
\end{notation}

\begin{obs}
  $\abs{B_S(g,r)} = \abs{B_S(r)} \leq \abs{S}^r$.
\end{obs}

\begin{definition}
  Sea $G$ grupo finitamente generado y sea $S \subset G$ un conjunto finito, simétrico, generador. Decimos que:
  \begin{enumerate}[(1)]
    \item $G$ tiene \wea{crecimiento exponencial} si $$\exists \lam>1, \quad \forall r \geq 1, \quad \abs{B_S(r)} \geq \lam^r.$$
    \item $G$ tiene \wea{crecimiento polinomial} si $$\exists p \in \R[x], \quad \forall r \geq 1, \quad \abs{B_S(r)} \leq p(r).$$
    \item $G$ tiene \wea{crecimiento intermedio} si no es ni polinomial ni exponencial.
  \end{enumerate}
\end{definition}

\begin{exercise}
  La definición anterior no depende de $S$.
\end{exercise}

\begin{example}
  Si $G$ es finito, entonces es de crecimiento polinomial. Se acota por un polinomio de grado el orden del grupo.
\end{example}

\begin{example}
  $\Z$ es de crecimiento polinomial. Si tomamos $S = \{-1,0,1\}$, entonces $\abs{B_S(r)} = 2r + 1$, tomamos $p(x) = 2x+1$.
\end{example}

\begin{example}
  $\Z^d$ es de crecimiento polinomial. Si elegimos
  \[S = \set{\sum_{i=1}^d a_ie_i \colon a_i \in \{-1,0,1\}}\]
  entonces $\abs{B_S(r)} = (2r+1)^d.$
\end{example}

\begin{definition}
  Decimos que un grupo $G$ es \wea{nilpotente} si al escribir
  \[G_0 = G, \quad G_{n+1} = \frac{G_n}{\Cent(G_n)}\]
  entonces la secuencia $\{G_n\}_n$ termina en $\{1_G\}$.
\end{definition}

\begin{exercise}
  El \wea{grupo de Heisenberg discreto} es
  \[H = \set{\matr{1 & x & z \\ 0 & 1 & y \\ 0 & 0 & 1} \in \GL_3(\Z) \colon x,y,z \in \Z} \leqslant \GL_3(\Z).\]
  $H$ es el ejemplo clásico de un grupo nilpotente infinito y finitamente generado. Tomamos
  \[S = \set{\matr{1&1&0\\0&1&0\\0&0&1}^\pm, \matr{1&0&1\\0&1&0\\0&0&1}^\pm, \matr{1&0&0\\0&1&1\\0&0&1}^\pm}.\]
  Pruebe que $H$ tiene crecimiento polinomial.
\end{exercise}

\begin{theorem}
  [Gromov]
  Un grupo $G$ finitamente generado tiene crecimiento polinomial si t solo si $\exists H \leqslant G$ con $[G \colon H] < \oo$ tal que $H$ es nilpotente.
\end{theorem}
\begin{proof}
  No veremos la demostración, ya que requiere muchas cosas complicadas.
\end{proof}

\newpage
\subsection{Grupos Libres}

\begin{definition}
  Sea $X$ un conjunto. Definimos
  \[X^\star = \bigcup_{n\geq0} X^n\]
  con la convención de que $X^0 = \{\e\}$.
\end{definition}

\begin{notation}
  $\omega = x_1 \dots x_n \in X^n \subset X^\star$ es una palabra de largo $n = \abs{\omega}$. i.e.: concatenación.
\end{notation}

\begin{example}
  Si $X = \{0,1\}$, entonces
  \[X^\star = \{\e,0,1,00,01,10,11,000,001,010,011, \dots\}.\]
\end{example}

\begin{definition}
  Sea $X$ un conjunto, considero a $X^-$ como un conjunto de ``inversos formales'':
  \[X^- = \{x^{-1} \colon x \in X\}.\]
\end{definition}

\begin{definition}
  Decimos que $(X,\op)$ es un \wea{monoide} si
  \begin{enumerate}[(1)]
    \item $X$ es un conjunto;
    \item $\op \colon X \times X \to X$ es una operación binaria;
    \item $\forall x,y,z \in X$, $\op(x,\op(y,z)) = \op(\op(x,y),z)$;
    \item $\exists 1 \in X$, $\forall x \in X$, $1x = x1 = x$.
  \end{enumerate}
\end{definition}

\begin{obs}
  Todos los grupos son monoides, pero los monoides no son necesariamente grupos, ya que les falta la existencia de inversos.
\end{obs}

\begin{obs}
  Consideramos $(X \cupdot X^{-1})^\star$ equipado con la concatenación:
  dados $u,v \in (X \cupdot X^{-1})^\star$ con $u = x_1 \dots x_n$ y $v = y_1 \dots y_m$, entonces
  \[u \c v = x_1 \dots x_n y_1 \dots y_m.\]
  Tenemos que $(X \cupdot X^{-1})^\star$ con la concatenación es un monoide.
\end{obs}

\begin{definition}
  Decimos que $u,v \in (X \cupdot X^{-1})^\star$ son \wea{similares} si $u$ puede obtenerse a partir de $v$ borrando o agregando un par de la forma $xx^{-1}$ o $x^{-1} x$ para algún $x \in X$.
\end{definition}

\begin{example}
  Si $X = \{a,b\}$, tenemos que $u = aba^{-1}ab$ y $v = abb$ son similares.
\end{example}

\begin{definition}
  Decimos que $u,v \in (X \cupdot X^{-1})^\star$ son \wea{equivalentes} si $\exists n \geq 1$ y $\w_0,\w_1,\dots,\w_n$ palabras tal que
  \[u = \w_0 v \w_1 v \w_2 \dots v \w_n = v.\]
  Se denota $u \sim v$.
\end{definition}

\begin{example}
  $abba^{-1}ab^{-1}b^{-1}abaa^{-1} \sim aab$.
\end{example}

\begin{obs}
  Ser palabras equivalentes es una relación de equivalencia.
\end{obs}

\begin{definition}
  Una palabra $\w \in (X \cupdot X^{-1})^\star$ se dice \wea{reducida} si no contiene copias de $xx^{-1}$ o $x^{-1}x$.
\end{definition}

\begin{obs}
  Cada clase de equivalencia contiene una única palabra reducida.
\end{obs}

\begin{obs}
  Dadas $u,v \in (X \cupdot X^{-1})^\star/\sim$, defino
  \[[u]_\sim \c [v]_\sim = [uv]_\sim.\]
  Esta operación está bien definida.
\end{obs}

\begin{example}
  Si $X = \{a,b\}$ y $u = abab^{-1} aa$, $v = a^{-1}a^{-1} bba$, entonces
  \begin{align*}
    uv &= a b a b^{-1} a a a^{-1} a^{-1} b b a \\
    &\sim a b a b^{-1} a a^{-1} b b a \\
    &\sim a b a b^{-1} b b a \\
    &\sim a b a b a
  \end{align*}
  y entonces $[uv]_\sim$ está representada por $ababa$.
\end{example}

\begin{definition}
  Sea $X$ un conjunto. El \wea{grupo libre generado por $X$} es
  \[F(X) = \frac{(X \cupdot X^{-1})^\star}{\sim}\]
  equipado con la concatenación.
\end{definition}

\begin{notation}
  Módulo isomorfías,
  \[F(X) = \{\w \in (X \cupdot X^{-1})^\star \text{ reducidas}\}.\]
\end{notation}

\begin{example}
  Si $X = \{a\}$, entonces
  \begin{align*}
    F(\{a\}) &= \{\w \in \{a,a^{-1}\}^\star \colon aa^{-1}, a^{-1} a \text{ no están en } \w\} \\
    &= \{a^n \colon n \geq 1\} \cup \{a^{-n} \colon n \geq 1\} \cup \{\e\} \\
    \simeq \Z.
  \end{align*}
\end{example}

\begin{example}
  Si $X = \{a,b\}$, entonces $F(X)$ son todas las palabras tales que no contiene ocurrencias de $aa^{-1}$, $a^{-1}a$, $bb^{-1}$, $b^{-1}b$. Consideremos $S = \{\e,a,b,a^{-1},b^{-1}\}$, entonces el grafo $\Cay(F(\{a,b\}),S)$ es
  \[\wea{Poner aca el grafo de Cayley de $F_2$.}\] %% no olvidar
\end{example}

\begin{exercise}
  Calcule explícitamente $B_S(r)$ y concluya que $F(\{a,b\})$ tiene crecimiento exponencial.
\end{exercise}

\begin{obs}
  Si $\abs X = \abs Y$, entonces $F(X) \simeq F(Y)$.
\end{obs}

\begin{notation}
  Dado $n \geq 1$ entero, denotamos $F_n$ al grupo libre generado por un conjunto de $n$ elementos.
\end{notation}

\fecha{Clase 9. Lunes 6 de Abril}

\begin{theorem}[\wea{Propiedad Universal de los Grupos Libres}]
  Sea $G$ un grupo, $X$ un conjunto y $f \colon X \to G$ función. Entonces $\exists \vphi \colon F(X) \to G$ homomorfismo tal que $\forall x \in X$, $\vphi(x) = f(x)$.

  Además, $F(X)$ es el único grupo que cumple esta propiedad.
\end{theorem}
\begin{proof}
  Clara Loh.
\end{proof}

\begin{corollary}
  Todo grupo es cociente de un grupo libre.
\end{corollary}
\begin{proof}
  $X = G$, $f = \id_G$.
\end{proof}

\begin{corollary}
  Si $G$ es grupo y $S \subset G$ es generador, entonces $G$ es cociente de $F(S)$.
\end{corollary}
\begin{proof}
  $X = S$, $f \colon S \to G$, $f(s) = s$.
\end{proof}

\begin{corollary}
  Si $G$ es generado por un conjunto $S$ finito de $n$ elementos, entonces $G$ es cociente de $F_n$.
  % Entraron a la casa del Luengo, se robaron todo,   uedó la pura zorra: la mamá del Luengo.
  %% Quiero dejar en claro que el señor Benjamín Maldonado Villalobos escribió esto, no yo (me está obligando a no borrarlo)
\end{corollary}

\newpage
\subsection{Presentaciones de Grupo}

\begin{definition}
  Sea $S$ un conjunto y sea $R \subset (S \cup S^{-1})^\star$. El \wea{grupo presentado por $S,R$} es
  \[\gen{S \mid R} = \frac{F(S)}{\ll R \gg}\]
  donde
  \[\ll R \gg = \bigcap_{\substack{N \normal F(S) \\ R \subset N}} N.\]
  Al conjunto $S$ se le denomina \wea{generadores}, y $R$ se le denomina \wea{relaciones}.
\end{definition}

\begin{obs}
  La intuición es que $\gen{S \mid R}$ es el grupo más grande generado por $S$ tal que toda palabra en $R$ es trivial.
\end{obs}

\begin{exercise}
  $\gen{S \mid \varnothing} \simeq F(S)$.
\end{exercise}

\begin{exercise}
  $\gen{\{a\} \mid \{a^n\}} \simeq \zmz n$.
\end{exercise}

\begin{exercise}
  $\gen{\{a,b\} \mid \{aba^{-1} b^{-1}\}} \simeq \Z^2$.
\end{exercise}

\begin{notation}
  \
  \begin{enumerate}[(1)]
    \item No escribiremos las llaves de los conjuntos $S$ y $R$. Si $S = \{s_1, \dots, s_n\}$ y $R = \{r_1, \dots, r_m\}$, entonces
    \[\gen{S \mid R} = \gen{s_1,\dots, s_n \mid r_1, \dots, r_m}.\]
    \item Las relaciones se pueden escribir con igualdad, como por ejemplo
    \[aba^{-1}b^{-1} \leadsto aba^{-1} b^{-1} = 1 \leadsto ab=ba\]
    \[\implies \gen{\{a,b\} \mid \{aba^{-1}b^{-1}\}} = \gen{a,b \mid ab=ba}.\]
  \end{enumerate}
\end{notation}

\begin{definition}
  Dado $S$ conjunto y $R \subset (S \cup S^{-1})^\star$, considere el monoide $(S \cup S^{-1})^\star$ y la relación $\sim_R$ dada por:
  $u \sim_R v$ si existe una secuencia finita
  \[u = u_0 \sim_R u_1 \sim_R \dots \sim_R u_n = v\]
  tal que $u_{i+1}$ se obtiene a partir de $u$ mediante:
  \begin{itemize}[(1)]
    \item Agregar o quitar un pat $ss^{-1}$ o $s^{-1}s$,
    \item Agregar o quitar $r$ o $r^{-1}$ con $r \in R$.
  \end{itemize}
\end{definition}

\begin{obs}
  $(S \cup S^{-1})^\star/\sim_R$ es un grupo con la concatenación.
\end{obs}

\begin{theorem}
  $$\frac{(S \cup S^{-1})^\star}{\sim_R} \simeq \frac{F(S)}{\ll R \gg}.$$
\end{theorem}

\begin{lemma}
  Si $u,v \in F(X)$ tal que $uv \in \ll R \gg$, entonces $\forall r \in (R \cup R^{-1})^\star$, $urv \in \ll R \gg$.
\end{lemma}

\begin{definition}
  Sea $S$ un conjunto de generadores y $S$ de relaciones. Decimos que
  \begin{itemize}[(1)]
    \item $\gen{S \mid R}$ es \wea{finitamente generado} si $\abs S < \oo$.
    \item $\gen{S \mid R}$ es \wea{finitamente presentado} si $\abs S < \oo$ y $\abs R < \oo$.
  \end{itemize}
  Más generalmente, un grupo $G$ se dice \wea{finitamente presentado} si existen $S$ y $R \subseteq (S \cup S^{-1})^\star$ finitos tales que
  \[G \simeq \gen{S \mid R}.\]
\end{definition}

\begin{example}
  $F_n, \Z^d$ y los grupos finitos son finitamente presentados.
\end{example}

\begin{example}
  Cualquier grupo dado por una presentación finita
  \[\mathrm{BS}(1,n) = \gen{a,b \mid ab=ba^n}.\]
\end{example}

\begin{example}
  $$\mathrm{FB}(n,k) = \gen{a_1,\dots,a_n \mid \forall \w \in \{a_1^\pm, \dots, a_n^\pm\}, \w^k = 1}.$$
\end{example}

\begin{obs}
  ¿Existe algún $n,k$ tal que $\mathrm{FB}(n,k)$ sea infinito? Este se conoce como el \wea{problema de Burnside} y se resolvió afirmativamente en los 40's.
\end{obs}

\begin{obs}
  Dada una presentación finita $S,R$, ¿se puede decidir algoritmicamente si $\gen{S \mid R} \simeq 1$? La respuesta es que no se puede. En particular, este problema es indecidible.
\end{obs}

\begin{definition}
  Dado $G = \gen{S}$ con $S$ finito. El \wea{problema de la palabra} de $G$ es
  \[\mathrm{WP}_S(G) = \{\w \in (S \cup S^{-1})^\star \colon \underline{\w}_G = 1_G\}.\]
\end{definition}

\begin{theorem}
  Si $G$ es un grupo finitamente presentado, el problema de decirdir si $\w \in (S \cup S^{-1})^\star$ pertenece a $\mathrm{WP}_S(G)$ es indecidible.
\end{theorem}

Sea $G$ un grupo. ¿Cómo probar que $F_n \embed G$? Hay una técnica bonita que permite hacer eso:

\begin{lemma}
  Sea $n \geq 2$, $X$ un conjunto y $a_1, \dots, a_n \in \Sym(X)$. Supongamos que existen subconjuntos de $X$ disjuntos a pares
  \[X_1^+, \dots, X_n^+, X_1^-, \dots, X_n^-\]
  tales que
  \[Z = \bigcup_{i=1}^n X_i^+ \cup \bigcup_{i=1}^n X_i^- \implies \forall i \cor{
    a_i (Z \setminus X_i^-) \subset X_i^+ \quad \land\quad a_i^{-1} (Z \setminus X_i^+) \subseteq X_i^-
  }\]
  entonces $\gen{a_1,\dots,a_n} \simeq F_n$.
\end{lemma}

\fecha{Clase 10. Miércoles 8 de Abril}

\begin{obs}
  Este lema nos permite incrustar $F_2$ en otro grupo.
\end{obs}

Una versión más elegante del lema del Ping Pong es la siguiente:

\begin{theorem}
  Sea $G \acts X$ una acción y sean $H_1,\dots, H_k \leqslant G$ (con $k \geq 2$) con al menos uno que no sea trivial.
  Supongamos que existen $X_1, \dots, X_k \subset X$ disjuntos de a pares tales que $\forall s \in \{1,\dots, k\}$, $\forall i \neq s$, $\forall h \in H_i \setminus \{1_G\}$ tenemos
  \[h(X_s) \subset X_i.\]
  Entonces
  \[\gen{\bigcup_{i=1}^k H_i}
  \simeq H_1 * H_2 * \dots * H_k
  \simeq \gen{\bigcup S_i \mid \bigcup R_i}\]
  donde $H_i = \gen{S_i \mid R_i}$ con los $S_i$ disjuntos de a pares.
\end{theorem}

\begin{theorem}
  $F_2 \embed \SL_2(\Z)$. Más precisamente,
  \[\gen{\matr{1 & 2 \\ 0 & 1}, \matr{1 & 0 \\ 2 & 1}} \simeq F_2.\]
\end{theorem}

\begin{example}
  Consideremos $X = \R P^1$ la \wea{recta proyectiva real}
  \[\R P^1 = \frac{\R^2 \setminus \{(0,0)\}}{x \sim \lam x}\]
  La acción de $\SL_2(\Z) \acts \R^2$ induce una acción $\SL_2(\Z) \acts \R P^1$. Sean $x,y \in \R^2 \setminus 0$ tal que $[x] = [y]$, entonces existe $\lam \neq 0$ tal que $x = \lam y$. Luego
  \[A \c x = A(\lam y) = \lam A y \implies Ax \sim Ay.\]
  Escribimos
  \[g_1 = \matr{1 & 2 \\ 0 & 1}, \quad g_2 = \matr{1 & 0 \\ 2 & 1}.\]
  Notemos que
  \[g_1 \matr{x \\ y} = \matr{x+2y \\ y}, \quad
  g_2 \matr{x \\ y} = \matr{x \\ 2x+y}.\]
  Consideramos
  \[H_1 = \gen{g_1} \simeq \Z, \quad
  H_2 = \gen{g_2} \simeq \Z.\]
  Consideremos
  \begin{align*}
    X_1 &= \{[x:y] \in \R P^1 \colon \abs x > \abs y\}, \\
    X_2 &= \{[x:y] \in \R P^1 \colon \abs y > \abs x\}.
  \end{align*}
  Notemos que si $[x:y] \in X_2$, entonces $g_1 \matr{x\\y} \in X_1$, y si $[x:y] \in X_1$, entonces $g_2 \matr{x\\y} \in X_2$.

  Si $h \in H_1$, entonces existe $n \in \Z$ tal que
  \[h = \matr{1 & 2n \\ 0 & 2}.\]
  Si $h \in H_1 \setminus \{1\}$, entonces $n \neq 0$. Notemos que
  \[h \matr{x\\y} = \matr{1 & 2n \\ 0 & 1} \matr{x\\y} = \matr{x + 2ny \\ y}.\]
  Luego si $[x:y] \in X_2$, tenemos que $h \c [x:y] \in X_1$. Similarmente si $[x:y] \in X_1$ y $\gamma \in H_2 \setminus \{1\}$, entonces $\gamma \c [x:y] \in X_2$.
  
  Por el lema del Ping Pong (elegante) tenemos que
  \[H_1 * H_2 \simeq \gen{\matr{1 & 2 \\ 0 & 1}, \matr{1 & 0 \\ 2 & 1}} \simeq F_2.\]
\end{example}

\begin{exercise}
  Pruebe que $F_n \embed F_2$ para $n \geq 2$.

  \emph{Hint:} Consideremos el conmutador $[F_2,F_2]$. Pruebe que $[F_2,F_2] \simeq F_\oo$, luego
  \[F_n \embed F_\oo \simeq [F_2,F_2] \leqslant F_2.\]
\end{exercise}

\newpage
\subsection{Introducción a lso Teoremas de Sylow}

Vamos a ver una parte del Teorema de clasificación de grupos abelianos finitamente generados.

\begin{theorem}
  Sea $G$ un grupo abeliano finitamente generado. Entonces existen $d \geq 0$, $p_1,\dots, p_k$ primos y $\alpha_1, \dots, \alpha_k \geq 1$ tales que
  \[G \simeq \Z^d \oplus \bigoplus_{i=1}^k \znz{p_i^{\alpha_i}}.\]
\end{theorem}

\begin{example}
  Supongamos que $G$ es $n$--generado, entonces $G$ es cociente de $\Z^n$. Es decir, existe $N \normal \Z^n$ tal que
  \[G \simeq \frac{Z^n}{N}.\]
  Notemos que $N = A \Z^n$ con $A$ matríz. Esta tiene descomposición de Smith:
  \[A = UST\]
  donde $U,T$ son matrices invertibles y $S$ es diagonal.
\end{example}

\begin{definition}
  Sea $G$ un grupo y $g \in G$. el \wea{orden de $g$} se define como
  \[\ord(g) = \min\{n \geq 2 \colon g^n = 1\},\]
  \[\ord(g) = \oo \iff \gen{g} \simeq \Z.\]
  $g$ se dice \wea{de torsión} si $\ord(g) < \oo$.
  $g$ se dice \wea{no de torsión} si $\ord(g) = \oo$.
\end{definition}

\begin{definition}
  Sea $G$ un grupo y $p$ primo. Decimos que $G$ es un \wea{$p$--grupo} si $\forall g \in G$, $\ord(g) = p^n$ para algún $n \geq 0$.
\end{definition}

\begin{obs}
  Si $G$ es finito, entonces $G$ es $p$--grupo si y solo si $\abs G = p^n$ para algún $n \geq 0$.

  Si $\abs G$ no es potencia de $p$, entonces existe $q \neq p$ primo tal que $q \mid \abs G$. Por el Teorema de Cauchy, $\zmz q \embed G$, y entonces $\exists a \in G$ tal que $\ord(a) = q$. $\to\gets$
\end{obs}

\begin{theorem}
  Todo grupo abeliano finito se descompone como suma de sus $p$--subgrupos maximales ($p$--subgrupos más grandes bajo contención).
\end{theorem}

\begin{definition}
  Dado un grupo $G$, definimos el \wea{exponente de $G$} como
  \[\exp(G) = \min\{n \geq 1 \colon \forall g \in G, g^n = 1\}.\]
\end{definition}

\begin{obs}
  $\exp(G) \leq \abs G$. Podría ser menor estricto:
  \[G = \znz2 \times \znz2\]
  tenemos que $\exp(G) = 2$ y $\abs G = 4$.
\end{obs}

\begin{obs}
  Si $\abs G < \oo$, entonces $\exp(G) \mid \abs G$.
\end{obs}

\newpage %% error mdframed
\begin{obs}
  Si $\abs G < \oo$, entonces
  \[\exp(G) = \lcm\{\ord(g) \colon g \in G\}.\]
\end{obs}
\begin{proof}
  Sea $M = \lcm\{\ord(g) \colon g \in G\}$. Notemos que $M \leq \exp(G)$, pues $\forall g \in G$, $\ord(g) \mid \exp(G)$.

  Por otro lado, $g^M = 1$ para todo $g \in G$, y justamente $\exp(G)$ es el menor entero positivo que cumple esto, por lo que $\exp(G) \leq M$. Concluimos que $\exp(G) = M$.
\end{proof}

\begin{notation}
  De aquí en adelante escribimos
  \[\abs G = p_1^{\a_1} \dots p_\ell^{\a_\ell}\]
  con $p_i$ primos y $\a_i \geq 1$ enteros.
  \[\exp(G) = p_1^{r_1} \dots p_\ell^{\a_\ell}\]
  con $r_i \geq 1$ enteros y $r_i \leq \a_i$.
\end{notation}

\begin{definition}
  Para $d \geq 1$ defino
  \begin{align*}
    E_d \olon G &\to G \\
    g &\mapsto g^d.
  \end{align*}
  Escribimos $G_d = \ker(E_d)$ y $E_d(G) = \im(E_d)$.
\end{definition}

\begin{proposition}
  \[1 \to G_d \to G \to E_d(G) \to 1\]
\end{proposition}

\begin{obs}
  $G_d \simeq 1 \iff \gcd(d,\abs G) = 1$.
\end{obs}

\begin{lemma}
  Sea $d = p_i^{r_i}$, entonces
  \begin{enumerate}[(1)]
    \item $p_i$ no divide a $\abs{E_d(G)}$.
    \item $G \simeq G_{p_i^{r_i}} \times E_{p_i^{r_i}}(G)$.
  \end{enumerate}
\end{lemma}

\newpage
\begin{obs}
  De acá sale el Teorema, pues inductivamente
  \[G \simeq G_{p_1^{r_1}} \times \dots \times G_{p_\ell^{r_\ell}}.\]
\end{obs}
\begin{proof}
  Supongamos que $p_i \mid \abs{E_d(G)}$. Entonces $\zmz{p_i} \embed E_d(G)$. Ñuego $\exists g \in E_d(G) \setminus \{1\}$ tal que $g^{p_i} =1$.
  Como $g \in E_d$, tenemos que $g = h^{p_i^{r_i}}$ para algún $h \in G \setminus \{1\}$.
  \[\implies 1 = g^{p_i} = h^{p_i^{r_i+1}}
  \implies \ord(h) = p_i^{r_i+1}.\]
  Esto contradice que $\exp(G) = p_1^{r_i} \dots p_\ell^{r_\ell}$.

  Por PTI, $$\frac{G}{G_d} \simeq E_d(G).$$
  En particular, $\abs G = \abs{G_d} \c \abs{E_d(G)}$. Luego $\abs{G_d} = p_i^{r_i}$ y $G_d \cap E_d(G) = \{1\}$.
  Como $G_d$ y $E_d(G)$ son normales en $G$, entonces
  \[G_d \c E_d(G) \simeq G_d \times E_d(G)\]
  Como $G = G_d \c E_d(G)$ se concluye (ejercicio).
\end{proof}

\fecha{Clase 11. Lunes 13 de Abril}

\newpage
\subsection{Productos Semidirectos}

\begin{definition}
  Sean $H,K$ dos grupos y sea $\a \colon K \to \Aut(H)$ un homomorfismo de grupos (i.e.: una acción $K \acts H$). El \wea{producto semidirecto} de $H$ y $K$ bajo $\a$, denotado por $H \rtimes_\a K$, es el conjunto $H \times K$ con la operación
  \[(h_1,k_1) \c (h_2,k_2) = (h_1 \c \a_{k_1}(h_2), k_1 \c k_2).\]
\end{definition}

\begin{notation}
  Escribimos $\a_k$ para el automorfismo $\a(k)$.
\end{notation}

\begin{theorem}
  $H \rtimes_\a K$ es grupo.
\end{theorem}

\begin{obs}\
  \begin{itemize}
    \item $H \rtimes_\a K$ contiene copias isomorfas de $H$ y $K$:
    \[ H \simeq H \times \{1_K\} \leqslant H \rtimes_\a K, \quad
    K \simeq \{1_H\} \times K \leqslant H \rtimes_\a K.\]
    \item $H \times \{1_K\} \cap \{1_H\} \times K = \{(1_H,1_K)\}.$
    \item $(1_H,k) \c (h,1_K) \c (1_H,k)^{-1} = (\a_k(h),1_K)$.
  \end{itemize}
  Luego $H \times \{1_K\}$ es normal en $H \rtimes_\a K$, luego la acción por conjugación de $\{1_H\} \times K$ sobre $H \times \{1_K\}$ es simplemente 
\end{obs}

Otra definición equivalente a la del producto semidirecto es la siguiente:

\begin{definition}
  Sea $G$ un grupo, y $H,K \leqslant G$ tales que
  \begin{itemize}[(i)]
    \item $HK = G$ (i.e.: $\forall g \in G, \exists h \in H, \exists k \in K, g = hk$).
    \item $H \cap K = \{1_G\}$.
    \item $H \normal G$.
  \end{itemize}
  Entonces diremos que $G$ es \wea{producto semidirecto} de $H$ y $K$.
\end{definition}

\begin{proposition}
  En el caso anterior, existe una acción natural $\a \colon K \o \Aut(H)$ con $\a_k(h) = khk^{-1}$, y se cumple que $G \simeq H \rtimes_\a K$, $hk \mapsto (h,k)$.
\end{proposition}

\begin{notation}
  Si $K = \Z$, la acción $\a \colon \Z \to \Aut(H)$ queda totalmente determinada por el automorfismo $\vphi = \a_1$. Por ende, escribimos $H \rtimes_\vphi \Z$ en vez de $H \rtimes_\a \Z$.

  En particular, si por ejemplo $H = \Z^d$, podemos escribir el producto como $\Z^d \rtimes_A \Z$ donde $A$ es la matriz asociada al automorfismo.
\end{notation}

\begin{definition}
  Dos acciones $\a,\beta \colon K \to \Aut(H)$ se dicen \wea{conjugadas} si existe un isomorfismo $\psi \colon H \to H$ tal que $\forall k \in K, \psi \circ \beta_k = \a_k \circ \psi$.
\end{definition}

\begin{proposition}
  Si $\a,\beta$ son conjugadas, entonces $H \rtimes_\a K \simeq H \rtimes_\beta K$.
\end{proposition}

\begin{proposition}
  Si $\a \colon K \to \Aut(H)$ es una acción y $\vphi \colon K \to K$ es un automorfismo, entonces $H \rtimes_\a K \simeq H \rtimes_{\a \circ \vphi} K$.
\end{proposition}

\renewcommand{\H}{\mathcal{H}}

\begin{example}
  Consideremos
  \[\H = \set{\matr{1 & a & c \\ 0 & 1 & b \\ 0 & 0 & 1} \colon a,b,c \in \Z}\]
  Notemos que la operación es
  \[\matr{1&a&b \\ 0&1&c \\ 0&0&1} \matr{1&x&y \\ 0&1&z \\ 0&0&1} = \matr{1 & a+x & az+b+y \\ 0 & 1 & cz \\ 0 & 0 & 1}\]
  Notemos que $\H \simeq \Z^2 \ltimes_A \Z$.
  \begin{align*}
    \Z^2 &\simeq \set{\matr{1&a&b \\ 0 & 1 & 0 \\ 0 & 0 & 1} \colon a,b \in \Z} = H, \\
    \Z &\simeq \set{\matr{1&0&0 \\ 0&1&c \\ 0&0&1} \colon c \in \Z} = K.
  \end{align*}
  \begin{itemize}
    \item $\H = HK$ %% Comprobar
    \item $H \cap K = \{I\}$ %% Comprobar
    \item $H,K \normal \H$ %% comprobar
  \end{itemize}
  Consideramos la acción $\Z \actsby A \Z^2$ dada por
  \[\vec x \mapsto A \vec x = \matr{1&0\\1&1} \vec x\]
  Luego $\H = \Z^2 \rtimes_A \Z$.
\end{example}

\begin{exercise}
  Considere $S_3 \rtimes \zmz2$, y sea $\sigma = (12)$. Muestre que hay un isomorfismo
  \[S_3 \rtimes \znz 2 \simeq S_3 \times \znz 2.\]
\end{exercise}

\subsection{Producto en Corona (Wreath Product)}

Supongamos que tenemos un conjunto de índices $I$ y una acción $K \acts I$ (i.e.: un homomorfismo $\a \colon K \o B_{ij}(I,I)$). Este $\a$ induce una acción del grupo $K$ en cualquier grupo $H^I$, que envía $(h_j)_{j \in I}$ a $(h_{\a(j)})_{j \in I}$. El producto semidirecto $H^I \rtimes K$ se llama \wea{producto en corona}.

\fecha{Clase 12. Miércoles 15 de Abril}

\begin{definition}
  Supongamos que $H,K$ son grupos cualesquiera, y $J$ es un conjunto de índices para el cual hay una acción $K \acts J$. Luego $$\bigoplus _{j\in J} H$$ es un grupo cuyos elementos son tuplas $(h_j)_{j \in J}$ de elementos de $H$ en las que solo finitas de ellas son distintas de $1_H$.
\end{definition}

\begin{obs}
  El grupo $K$ actúa sobre $\bigoplus_{j \in J} H$ permutando los elementos de una tupla según sus índices:
  \[k \c (h_j)_{j \in J} \mapsto (h_{k \c j})_{j \in J}.\]
  Esta se conoce como la \wea{representación de Koopman}.
\end{obs}

\begin{definition}
  El producto semidirecto $\ds \bigoplus_{j \in J} H \rtimes K$ se denomina \wea{producto en corona} de $H$ y $K$, y se denota $H \wr K$.
\end{definition}

\begin{obs}
  La restricción de ``soporte finito'' a veces no se considera en la definición.
\end{obs}

\begin{obs}
  En general, si $\abs J = \oo$, $\ds \bigoplus_{j \in J} H$ no es finitamente generado.
  Sin embargo, si $H$ y $K$ son finitamente generados, entonces $H \wr K$ también lo es.
\end{obs}

\fecha{Clase 13. Lunes 20 de Abril}

\begin{definition}
  Sean $H$ y $K$ dos grupos. Sea $I$ un conjunto de índices y $K \actsby\a I$ acción de grupos.
  \begin{itemize}
    \item $\ds \bigoplus_{i \in I} H \coloneqq \{(h_i)_{i \in I} \colon h_i \in H \land \exists F \subseteq I, \forall i \in I \setminus F, h_i = 1_H\}.$ Este es un grupo con la operación elemento a elemento.
    \item \wea{Representación de Koopman:} Acción $\ds \tilde\a \colon K \acts \bigoplus_{i \in I} H$ dada por $k \c (h_i)_{i \in I} = (h_{k^{-1} \c i})_{i \in I}$.
    \item \wea{Producto en Corona (Wreath Product):} $\ds H \wr_\a K \coloneqq \bigoplus_{i \in I} H \rtimes_{\tilde\a} K$.
  \end{itemize}
\end{definition}

Usualmente $I = K$ y la acción $\a$ es multiplicación por la izquierda, en cuyo caso omitimos el subíndice.

\begin{example}
  $W_d = \{\text{isometrías de $[0,1]^d$}\} \simeq \zmz2 \wr S_d$.
\end{example}

\begin{example}
  [\wea{Lamplighter}]
  $L = \zmz 2 \wr \Z$. Este es un grupo finitamente generado con un subgrupo no finitamente generado ``grande''.
  \[L = \gen{\{s_i\}_{i \in \Z} \cup \{t\} \colon s_i^2 = 1, s_i t = ts_{i+1}} = \gen{s_0, t}.\]
\end{example}

\begin{obs}
  Si $H$ y $K$ son finitamente generados, entonces $H \wr K$ también es finitamente generado.
\end{obs}

\newpage
\subsection{Teoremas de Sylow}

\begin{theorem}
  [\wea{Estructura de grupos abelianos}]
  Si $A$ es abeliano y finitamente generado, entonces $\exists r,d_1,\dots, d_n \in \N$ tales que $d_i \mid d_{i+1}$ y
  \[A \simeq \Z^r \times \znz{d_1} \times \dots \times \znz{d_n}.\]
\end{theorem}

\begin{definition}
  Un \wea{$p$--grupo} es un grupo $G$ en el que todos sus elementos tienen orden una potencia de $p$.
\end{definition}

\begin{example}
  $\zmz{p^n}$ es un $p$--grupo. $\zmz2 \times \zmz2$ es un 2--grupo.
\end{example}

\begin{proposition}
  Un $p$--grupo finito tiene orden $p^k$, y  recíprocamente si $\abs G = p^k$, entonces $G$ es un $p$--grupo.
\end{proposition}

\begin{theorem}
  [\wea{Cauchy}]
  Si $p$ es primo y divide a $\abs G < \oo$, entonces $G$ tiene un elemento de orden $p$.
\end{theorem}

El objetivo de hoy es encontrar una especie de recíproco para el teorema de Lagrange.

\begin{lemma}
  Sea $G$ un $p$--grupo finito y $G \actsby\a X$ una acción. Entonces
  \[\abs X \equiv \abs{\Fix_G(X)} \pmod p\]
  donde $\Fix_G(X) = \{x \in X \colon \forall g \in G, g \c x = x\}.$
\end{lemma}
\begin{proof}
  Sea $X/G = \{\orb(x) \colon x \in X\}$ el espacio de órbitas. Luego $x \in \Fix_G(X) \iff \{x\} \in X/G$.
  Por el teorema órbita--estabilizador, para todo $x \in X$ tenemos
  \[\abs{\orb(x)} = \frac{\abs G}{\abs{\stab(x)}},\]
  por lo que si $x$ no es un punto fijo, $\abs{\orb(x)}$ es una potencia de $p$, y luego $\abs{\orb(x)} \equiv 0 \pmod p$. Notemos que
  \[X = \bigcupdot_{\Omega \in X/G} \Omega\]
  \begin{align*}
    \implies \abs X &= \sum_x \abs{\orb(x)} \\
    &= \sum_{x \in \Fix_G(X)} 1 + \sum_{x \text{ no fijo}} \abs{\orb(x)} \\
    &= \abs{\Fix_G(X)} + \sum_{x \text{ no fijo}} \underbrace{\abs{\orb(x)}}_{\equiv 0 \mod p} \\
    &\equiv \abs{Fix_G(X)} \pmod p.
  \end{align*}
\end{proof}

\begin{theorem}
  [\wea{Primer Teorema de Sylow}]
  Sea $G$ un grupo finito y supongamos que $\abs G = p^k m$ con $p \nmid m$ primo. Entonces $G$ posee un subgrupo de orden $p^k$.
\end{theorem}

\begin{definition}
  Sea $G$ un grupo finito con $\abs G = p^k m$, donde $p \nmid m$ es primo. Un \wea{$p$--subgrupo de Sylow} de $G$ es un subgrupo de orden $p^k$.
\end{definition}

\begin{theorem}
  [\wea{Segundo Teorema de Sylow}]
  Si $H_1,H_2$ son $p$--subgrupos de Sylow, entonces son conugados (y en particular, isomorfos).
\end{theorem}

\begin{theorem}
  [\wea{Tercer Teorema de Sylow}]
  Sea $n_p$ el número de $p$--subgrupos de Sylow de un grupo finito $G$ (con $\abs G = p^k m$, $p \nmid m$).
  Entonces
  \begin{enumerate}[(i)]
    \item $n_p \equiv 1 \pmod p$
    \item $n_p \mid m$
    \item Si $H$ es un $p$--subgrupo de Sylow de $G$, entonces $n_p = [G \colon \Norm_G(H)]$.
  \end{enumerate}
\end{theorem}

\begin{example}
  El siguiente subgrupo de $S_4$ es un grupo de orden 12 sin subgrupos de orden 6:
  \[\{\id, (21)(43), (13)(24), \dots\} = A_4\]
  Más generalmente, para $n \geq 5$, $A_n$ es un grupo de orden $n!/2$ que no puede poseer subgrupos de orden $n!/4$ ya que serían subgrupos normales, y $A_n$ es simple.
\end{example}

\newpage
\fecha{Clase 14. Miércoles 22 de Abril}
\subsection{Paradoja de Banach--Tarski}

\begin{definition}
  Sea $G \acts X$ una acción. Diremos que $A, B \subset X$ son \wea{$G$--equidescomponibles} si $\exists A_1, \dots, A_n, B_1, \dots, B_n \subset X$ y $g_1, \dots, g_n \in G$ tales que 
  \[ A = \bigcupdot_{i=1}^n A_i, \quad B = \bigcupdot_{i=1}^n B_i, \quad B_i = g_i A_i, \quad \forall i. \]
  Se denota $A \sim_{\mathrm{eq}} B$.
\end{definition}

\begin{example}
  $G = \R^2$,  $X = \R^2$, $G \acts \R^2$ por traslación. Si tomamos $A = [-1,1)^2$ y $B = [0,2) \times [0,1)$, tenemos la descomposición 
  \[ A_1 = [-1,1) \times [-1,0), \quad A_2 = [-1,1) \times [0,1), \quad B_1 = [0,1) \times [0,1), \quad [1,2) \times [0,1) \]
  entonces tenemos que $B_1 = p_1 + A_1, \quad B_2 = p_2 + A_2$. %% la cague aca, arreglar los conjuntos. deberia cortar $A$ en dos trozos, y poner uno al lado del otro.
\end{example}

\begin{obs}
  $\sim_\mathrm{eq}$ es relación de equivalencia.
\end{obs}

\begin{theorem}
  [Banach--Tarski]
  Sea $G = \R^3 \rtimes \mathrm{SO}_3(\R)$ y $G \acts \R^3$ la acción natural. Si $A,B \subset \R^3$ son acotados de interior no vacío, entonces $A \sim_{\mathrm{eq}} B$.
\end{theorem}

\begin{obs}
  En particular $B \sim_\mathrm{eq} B_1(0) \cup ((3,0) + B_1(0))$.
\end{obs}

\begin{definition}
  Sea $G \acts X$. Un conjunto $Y \subset X$ se duce \wea{$G$--paradojal} si $\exists A_1, \dots, A_n \subset X$, $B_1, \dots, B_m \subset X$, $g_1, \dots, g_n \in G$, $h_1, \dots, h_m \in G$ tales que
  \[ Y = \bigcupdot_{i=1}^n A_i \cupdot \bigcupdot_{j = 1}^m B_j = \bigcupdot_{i=1}^n g_i A_i  = \bigcupdot_{j=1}^m h_j B_j.\]
\end{definition}

\begin{definition}
  Un grupo $G$ se dice \wea{paradojal} si $G \acts G$ por traslación izquierda es $G$--paradojal.
\end{definition}

\begin{theorem}
  $F_2$ es paradojal.
\end{theorem}
\begin{proof}
  Tomamos $A_1$ como el conjunto de todas las palabras que empiezan con $a^{-1}$ y $A_2$ el conjunto de palabras que empiezan con $a$. Sea $B_1$ el conjunto de palabras que empiezan con $b^{-1}$ o que son potencia de $b$. Sea $B_2 = F_2 \setminus (A_1 \cup A_2 \cup B_1)$. Luego $F_2 = A_1 \cup A_2 \cup B_1 \cup B_2$, $g_1 = a$, $g_2 = 1$, $h_1 = b$, $h_2 = 1$. Entonces 
  \[ F_2 = aA_1 \cupdot A_2 = bB_1 \cupdot B_2. \]
\end{proof}

\begin{proposition}
  Sea $G$ un grupo paradojal y $G \acts X$ una acción libre. Entonces $X$ es $G$--paradojal.
\end{proposition}
\begin{proof}
  Como $G$ es paradojal, existen $A_1, \dots, A_n, B_1, \dots, B_m, g_1, \dots, g_n, h_1, \dots, h_m$ que cumplen la propiedad de $G$ paradojal.
  Sea $R \subset X$ tal que $\forall x \in X$ 
  \[ \abs{\orb_G(x) \cap R} = 1 \]
  (Esto se puede por axioma de elección)
  Luego definimos $\forall i, \forall j$, 
  \[ A_i' = A_i \c R, \quad B_j' = B_j \c R \]
  Notemos que 
  \[ X = G \c R = \bigcupdot_{i=1}^n A_i' \cupdot \bigcupdot_{j=1}^m B_j'
  = \bigcup_{i=1}^n g_i A_i' = \bigcup_{j=1}^m h_j B_j' \]
  donde $\bigcupdot_{i=1}^n A_i'$ y $\bigcupdot_{j=1}^m B_j'$ son disjuntos, ya que la acción es libre.
\end{proof}

\begin{corollary}
  Si $H \embed G$ y $H$ es paradojal, entonces $G$ es paradojal.
\end{corollary}

\begin{corollary}
  $\mathrm{SO}_n(\R)$ es paradojal para $n \geq 3$.
\end{corollary}
\begin{proof}
  En el control 2 mostramos que $F_2 \embed \mathrm{SO}_3(\R) \embed \mathrm{SO}_n(\R)$ para $n \geq 3$. Eventualmente añadiré la demostración acá. %% no olvidar
\end{proof}

\begin{definition}
  $S^2 = \{x \in \R^3 \colon \norm{x}_2 = 1\}$.
\end{definition}

\begin{proposition}
  $\exists D \subset S^2$ contable tal que $S^2 \setminus D$ es $\mathrm{SO}_3(\R)$--paradojal.
\end{proposition}
\begin{proof}
  Recordemos que $F_2 \embed \mathrm{SO}_3(\R)$. Sean $A,B \in \mathrm{SO}_3(\R)$ tales que 
  \[ F_2 \simeq \gen{A,B} \leqslant \mathrm{SO}_3(\R). \]
  Defino 
  \[ D = \{x \in S^2 \colon \exists g \in \gen{A,B} \setminus \{1\}, g \c x = x\} \]
  (Ejercicio: mostrar que $D$ es contable) %% no olvidar.
  Consideremos el conjunto 
  \[ \tilde D = \{z \in S^2 \colon \exists x \in D, \exists g \in \gen{A,B}, z = g \c x\}. \]
  $\tilde D$ es contable. Luego $S^2 \setminus \tilde D$ es $\gen{A,B}$--invariante. $F_2 \simeq \gen{A,B} \acts S^2 \setminus \tilde D$ es libre, luego $\gen{A,B} \acts S^2 \setminus \tilde D$ es paradojal.
  Por lo tanto $S^2 \setminus \tilde D$ es $\gen{A,B}$--paradojal, luego $S^2 \setminus \tilde D$ es $\mathrm{SO}_3(\mathbb{R})$--paradojal.
\end{proof}

\begin{proposition}
  $S^2 \sim_\mathrm{eq} S^2 \setminus D$ para todo $D$ contable.
\end{proposition}
\begin{proof}
  Como $D$ es contable, existe una recta $L \subset \R^3$ que pasa por 0 y $L \cap D = \varnothing$. Definimos 
    \[ W = \left\{ \theta \in [0,2\pi) \colon \exists x \in D, \exists n \in \Z, \rho_\theta^n x \in D \right\} \]
    donde $\rho_\theta$ es la rotación en ángulo $\theta$ en torno a $L$.

    El conjunto $W$ es contable, luego existe \( \theta \in [0,2\pi) \setminus W  \). Como $\rho_\theta \in \mathrm{SO}_3(\R)$, tenemos que $\forall n \in \Z \setminus \{0\}$, $\rho_\theta^n D \cap D = \varnothing$. Definimos 
    \[ \tilde D = \bigcup_{n \geq 0} \rho^n D \]
    Luego 
    \[ S^2 = \tilde D \cupdot (S^2 \setminus \tilde D) \sim_\mathrm{eq} \rho \tilde D \cupdot (S^2 \setminus \tilde D) = S^2 \setminus D \]
    pues $\rho \tilde D = \bigcup_{n \geq 1} \rho^n D$.
\end{proof}

\begin{proposition}
  Sea $G \acts X$ y $A,B \subset X$ tales que $A \sim_\mathrm{eq} B$. Entonces $A$ es $G$--paradojal $\iff B$ es $G$--paradojal.
\end{proposition}

\begin{corollary}
  $\mathrm{SO}_3(\R) \acts \R^3 \implies S^2$ es $\mathrm{SO}_3(\R)$--paradojal.
\end{corollary}

\begin{theorem}
  $\forall r > 0$, $B_r(0)$ es $\R^3 \rtimes \mathrm{SO}_3(\mathbb{R})$--paradojal.
\end{theorem}

\begin{obs}
  Como $S^2$ es $\mathrm{SO}_3(\mathbb{R})$--paradojal, entonces $B_r(0) \setminus \{0\}$ es $\mathrm{SO}_3$--paradojal.

  Para $0$, sea $\rho \in \R^3 \rtimes \mathrm{SO}_3$ una rotación de ángulo irracional en torno a la recta $\R \times \{1/2\} \times \{1/2\}$. Defino $D = \{\rho^n(0) \colon n \geq 0\}$, $\rho D = D \setminus \{0\}$.
  Luego 
  \[ B_r(0) = D \cup B_r(0) \setminus D \sim_\mathrm{eq} \rho D \cup B_r(0) \setminus D = B_r(0) \setminus \{0\} \]
  Por lo tanto $B_r(0)$ es $\R^2 \rtimes \mathrm{SO}_3$--paradojal.
\end{obs}

\begin{definition}
  Sea $G \acts X$ y $A,B \subset X$- Decimos que $A \leq_\mathrm{eq} B$ (\wea{subequidescomponible}) si $\exists C \subset B$ tal que $A \sim_\mathrm{eq} C$.
\end{definition}

\begin{proposition}
  $(A \leq_\mathrm{eq} B \land B \leq_\mathrm{eq} A) \implies A \sim_\mathrm{eq} B$.
\end{proposition}
\begin{proof}
  Sea $B_1 \subset B$ y $A_1 \subset A$ tales que $A \sim_\mathrm{eq} B_1$ y $B \sim_\mathrm{eq} A_1$. Sean $\varphi \colon A \to B_1$ y $\psi \colon B \to A_1$ biyecciones. Definimos $C_0 = A \setminus A_1$, y para todo $n \geq 0$ definimos $C_{n+1} = \psi(\vphi(C_n))$.

  Definimos $C = \bigcup_{n \geq 0} C_n$. Luego $\psi^{-1}(A \setminus C) = B \setminus \vphi(C)$. De acá se puede mostrar que $A \setminus C \sim_\mathrm{eq} B \setminus \vphi(C)$ y $C \sim_\mathrm{eq} \vphi(C)$ (Ejercicio). %% no olvidar
  Por lo tanto 
  \[ A = A \setminus C \cupdot C \sim_\mathrm{eq} B \setminus \vphi(C) \cupdot \vphi(C) = B. \]
\end{proof}

Ahora demostraremos Banach--Tarski:
\begin{proof}
  Por la proposición anterior, basta probar que $A \leq_\mathrm{eq} B$. Como $A$ es acotado, $\exists r > 0$ tal que $A \subset B_r(0)$. Sea $x \in \mathrm{int}(B) \neq \varnothing$. Por lo tanto $\exists \varepsilon>0$ tal que $B_\varepsilon(x) \subset B$. Como $B_r(0)$ es compacto, entonces existen $x_1, \dots, x_n$ tales que 
  \[ B_r(0) \subset \bigcup_{i=1}^n (x_i + B_\varepsilon(0)). \]
  Sean $y_1, \dots, y_n \in \R^3$ tales que $(y_i + B_\e(x))$ son disjuntos de a pares. Luego 
  \[ A \subset B_r(0) \subset \bigcup_{i=1}^n (x_i + B_\e(x))
  \leq_\mathrm{eq} \bigcupdot_{i=1}^n (y_i + B_\e(x))
  \sim_\mathrm{eq} B_\e(x) \subset B. \]
  Por lo tanto $A \leq_\mathrm{eq} B$.
\end{proof}

\begin{obs}
  En cualquier descomposición paradojal de $B_r(0)$ en $\R^3$ debe ser al menos un conjunto que no es Boreliano.
\end{obs}

\begin{definition}
  $G$ es \wea{promediable} si $\forall K \subset G$ finito y $\forall \delta>0$, $\exists F \subset G$ finito tal que 
  \[ \abs{FK \setminus F} \leq \delta \abs{F}. \]
\end{definition}

\begin{obs}
  $G$ es paradojal si y solo si no es promediable. i.e.: $\exists K \subset G$ finito y $\exists \delta > 0$ tales que $\forall F \subset G$ finito 
  \[ \abs{FK \setminus F} \geq \delta \abs F. \]
\end{obs}

%\fecha{Clase 15.}
%\fecha{Clase 16.}
%\fecha{Clase 17.}
%\fecha{Clase 18.}
%\fecha{Clase 19.}
%\fecha{Clase 20.}
%\fecha{Clase 21.}
%\fecha{Clase 22.}
%\fecha{Clase 23.}
%\fecha{Clase 24.}
%\fecha{Clase 25.}
%\fecha{Clase 26.}
%\fecha{Clase 27.}
%\fecha{Clase 28.}
%\fecha{Clase 29.}
%\fecha{Clase 30.}


\end{document}